\documentclass[11pt]{article}

\usepackage[margin=1.1in]{geometry}
\usepackage{amsmath,amssymb,amsthm,mathtools}
\usepackage{enumitem}
\usepackage{hyperref}
\usepackage{microtype}
\usepackage{booktabs}
\usepackage{mathrsfs}

\hypersetup{
  colorlinks=true,
  linkcolor=blue,
  citecolor=blue,
  urlcolor=blue,
  hypertexnames=false,
  pdftitle={Stratified Hammersley Laws for Luce Permutations},
  pdfauthor={Christopher D. Long}
}

\newtheorem{theorem}{Theorem}[section]
\newtheorem{proposition}[theorem]{Proposition}
\newtheorem{lemma}[theorem]{Lemma}
\newtheorem{corollary}[theorem]{Corollary}
\newtheorem{conjecture}[theorem]{Conjecture}
\theoremstyle{remark}
\newtheorem{remark}[theorem]{Remark}
\newtheorem{definition}[theorem]{Definition}
\newtheorem{warning}[theorem]{Warning}

\newcommand{\Pp}{\mathbb P}
\newcommand{\Ee}{\mathbb E}
\newcommand{\Rr}{\mathbb R}
\newcommand{\Nn}{\mathbb N}
\newcommand{\one}{\mathbf 1}
\newcommand{\Exp}{\operatorname{Exp}}
\newcommand{\LIS}{\operatorname{LIS}}
\newcommand{\LDS}{\operatorname{LDS}}
\newcommand{\Var}{\operatorname{Var}}
\newcommand{\dd}{\,d}
\newcommand{\eps}{\varepsilon}
\newcommand{\partialmark}{\dagger}
\newcommand{\calA}{\mathcal A}
\newcommand{\calL}{\mathcal L}
\newcommand{\calU}{\mathcal U}
\newcommand{\calV}{\mathcal V}
\newcommand{\norm}[1]{\left\lVert #1\right\rVert}
\newcommand{\abs}[1]{\left|#1\right|}
\newcommand{\dist}{\operatorname{dist}}

\title{Stratified Hammersley Laws for Luce Permutations}
\author{Christopher D. Long\thanks{\texttt{galizur@gmail.com}}}
\date{Draft --- October 2026}


\begin{document}
\maketitle

\begin{abstract}
We prove a Hammersley shape theorem for independent point arrays with exactly
one point opportunity on each deterministic vertical fibre.  If the density
in row $i$ converges uniformly to $q(i/n,\cdot)$ for a continuous profile
$q$, then the complete
longest-chain field, divided by $\sqrt n$, converges uniformly to the
inhomogeneous Hammersley value function
\[
  \calU_q(x,y)
  =2\sup_{\gamma:(0,0)\nearrow(x,y)}
  \int_0^1\sqrt{q(\gamma_1(t),\gamma_2(t))
  \gamma_1'(t)\gamma_2'(t)}\,\dd t.
\]
The convergence holds with fixed-deviation probability
$O(e^{-c\sqrt n})$, in every finite $L^p$, and almost surely under arbitrary
couplings.  We also prove localization of all near-maximal chains around the
compact set of maximizing curves.

The stratified theorem yields a general shape law for rankings generated by
independent, nonidentically distributed scores.  This includes Gaussian
Thurstone--Mosteller rankings and, after a bounded compactification of
exponential clocks, Luce permutations.  For every positive continuous Luce
profile we obtain the full permutation shape field, longest increasing and
decreasing subsequence laws, robustness under a sublinear set of bounded-rate
errors, and stability in the profile.  These results establish first-order
LIS asymptotics for the stated classes, addressing the question posed by
Chatterjee, Diaconis, and Kim and by Borga, Chatterjee, and Diaconis.

Finally, we develop a dyadic corner estimate for endpoint-degenerate rate
profiles.  It proves that the normalization remains $\sqrt n$ under dyadic
endpoint comparability, in particular for all regularly varying sampled
profiles.  As a canonical consequence, Luce permutations with Sukhatme
weights $\theta_{n,i}=n-i+1$ satisfy singular increasing- and
decreasing-subsequence laws.  Reversing the labels gives the corresponding
laws for $\theta_{n,i}=i$, with the two variational constants exchanged.
\end{abstract}

\tableofcontents

\section{Introduction}

The Hammersley process and the longest increasing subsequence of a uniform
random permutation are basic models of planar random growth.  The limit
\[
        \frac{\LIS(\pi_n)}{\sqrt n}\longrightarrow 2
\]
was obtained by Logan--Shepp \cite{LoganShepp1977} and, independently,
by Vershik--Kerov \cite{VershikKerov1977}.  See also the
Hammersley-process formulation of Aldous--Diaconis
\cite{AldousDiaconis1995}.  For independent
planar samples with a nonconstant density, Deuschel and Zeitouni proved the
corresponding variational law and localization of optimal chains
\cite{DeuschelZeitouni1995}, and later determined the fixed-deviation
large-deviation scales \cite{DeuschelZeitouni1999}.  Related permuton and RSK
limits were developed by Sj\"ostrand \cite{Sjostrand2023} and Dubach
\cite{Dubach2023,Dubach2024}.

For Mallows permutations $\pi_n$, with probability proportional to
$q_n^{\operatorname{inv}(\pi_n)}$, Mueller and Starr proved a $\sqrt n$-scale
LIS law when $n(1-q_n)\to\beta\in\mathbb R$
\cite{MuellerStarr2013}.  Bhatnagar and Peled proved
$\LIS(\pi_n)/(n\sqrt{1-q_n})\to1$ when $q_n\to1$ and
$n(1-q_n)\to\infty$ \cite{BhatnagarPeled2015}.

The Luce model \cite{Luce1959} is a natural nonuniform distribution on
permutations, obtained by sampling labels without replacement with
probabilities proportional to prescribed positive weights.  Equivalently, one ranks independent exponential
clocks with the corresponding rates.  Chatterjee, Diaconis, and Kim asked for
the distribution of the longest increasing subsequence
\cite[Introduction]{ChatterjeeDiaconisKim2024}.  Borga, Chatterjee, and
Diaconis subsequently established the Luce permuton and local limits and
again singled out the LIS problem in their Section~5.3
\cite{BorgaChatterjeeDiaconis2026}.  The two papers use inverse permutation
conventions, but $\LIS(\sigma)=\LIS(\sigma^{-1})$, so the distinction is
irrelevant here.

A permuton limit by itself cannot determine the $\sqrt n$-scale LIS: a
vanishing fraction of planted points can preserve the permuton while creating
a much longer monotone subsequence.  The clock representation supplies the
missing microscopic structure.  If $T_{n,i}$ is the clock attached to label
$i$, then exactly
\[
 \LIS(\sigma_n)
 =\max\{k:i_1<\cdots<i_k,\ T_{n,i_1}<\cdots<T_{n,i_k}\}.
\]
Thus the terminal LIS is already a longest-chain problem for independent
points.  The resulting cloud, however, has one point on every deterministic
vertical fibre; it is neither an i.i.d. planar sample nor a Poisson cloud.
For $n$ independent uniform points in the unit square, the probability of
placing exactly one point in each of $n$ equal vertical strips is
\[
        \frac{n!}{n^n}\sim\sqrt{2\pi n}\,e^{-n}.
\]
For a homogeneous Poisson cloud of intensity $n$ on the square, the same
event has probability $e^{-n}$.  Conditioning on this event therefore loses
a factor $e^{n+o(n)}$, too large to transfer the $e^{-c\sqrt n}$ deviation
bounds used here.  The first purpose of this paper is to prove the
appropriate stratified Hammersley theorem directly.

The proof uses a common-core decomposition.  Within a thin vertical strip,
every fibre law contains a common subdensity.  Conditional on selecting
this core, the vertical ranks are an exact uniform permutation.  The residual
subdensities are dominated by a sparse uniform proposal cloud.  Concatenating
the local chains produces a discrete variational problem, while a finite
vertical mesh controls the random block exit heights.  This gives the full
shape field and quantitative concentration.  A restricted version of the
same block upper bound proves localization near the maximizing curves.

The second purpose is to isolate the general ranking principle.  The reduction
does not depend on exponentials: it applies to independent continuous scores
whose transformed densities approach a continuous profile after one common
increasing compactification.
We formulate the induced permuton density and prove the complete shape law.
Gaussian random-utility rankings of Thurstone--Mosteller type are an immediate
example \cite{Thurstone1927,Mosteller1951}.  The Luce theorem is obtained by
specializing this independent-score result to exponential clocks.

The third purpose is to treat vanishing endpoint profiles.  A single bounded
compactification no longer works when the smallest rate tends to zero.  We
instead divide the terminal labels into dyadic rate-location blocks.  If the
rates are comparable inside every such block, the chain supported by the last
$m$ labels has exponential moments on the $\sqrt m$ scale.  Truncating the
endpoint and then sending the truncation to zero proves an ordinary
$\sqrt n$ Hammersley law.  Potter bounds show that sampled regularly varying
profiles satisfy the dyadic hypothesis \cite{BinghamGoldieTeugels1987}.  In
particular, this gives first-order increasing- and decreasing-subsequence laws
for the canonical Sukhatme weights $\theta_{n,i}=n-i+1$ considered in both
Luce papers, and for their label reversal.  Fluctuation questions are discussed
in Section~\ref{sec:further}.

\subsection{The regular-profile Luce theorem}

Let $f\in C[0,1]$ satisfy
\[
        0<m_f\le f(x)\le M_f<\infty.
\]
For $t\ge0$, define
\begin{equation}\label{eq:A-B-def}
        A_f(t)=\int_0^1(1-e^{-f(u)t})\,\dd u,
        \qquad
        B_f(t)=A_f'(t)=\int_0^1 f(u)e^{-f(u)t}\,\dd u.
\end{equation}
Then $A_f$ is a strictly increasing bijection from $[0,\infty)$ to $[0,1)$.
Write $\tau_f=A_f^{-1}$ and define
\begin{equation}\label{eq:rho-def-intro}
        \rho_f(x,y)
        =\frac{f(x)e^{-f(x)\tau_f(y)}}{B_f(\tau_f(y))},
        \qquad 0\le x\le1,\quad 0<y<1.
\end{equation}
This is the Luce permuton density identified by Borga, Chatterjee, and
Diaconis \cite[Proposition~2.2]{BorgaChatterjeeDiaconis2026}.

Let $\calA_+(x,y)$ be the set of absolutely continuous coordinatewise
nondecreasing paths from $(0,0)$ to $(x,y)$, and put
\begin{equation}\label{eq:Lrho-intro}
 \calU_{\rho_f}(x,y)
 =2\sup_{\gamma\in\calA_+(x,y)}
   \int_0^1
   \sqrt{\rho_f(\gamma_1(t),\gamma_2(t))
          \gamma_1'(t)\gamma_2'(t)}\,\dd t,
\end{equation}
with endpoint integrals interpreted by monotone truncation when necessary.

A Luce permutation with rates $\theta_{n,1},\dots,\theta_{n,n}$ is generated
by independent clocks
\[
        T_{n,i}\sim\Exp(\theta_{n,i}),
        \qquad
        \sigma_n(i)=\sum_{j=1}^n\one_{\{T_{n,j}\le T_{n,i}\}}.
\]
For $x,y\in[0,1]$, let $\Lambda_n(x,y)$ be the largest $k$ for which there
exist
\[
        1\le i_1<\cdots<i_k\le\lfloor nx\rfloor,
        \qquad
        \sigma_n(i_1)<\cdots<\sigma_n(i_k)\le\lfloor ny\rfloor.
\]
Thus $\Lambda_n(1,1)=\LIS(\sigma_n)$.

\begin{theorem}[Regular-profile Luce shape theorem]\label{thm:main-luce}
Suppose
\begin{equation}\label{eq:uniform-profile-hyp}
        \theta_{n,i}=a_n\bigl(f(i/n)+\eps_{n,i}\bigr),
        \qquad
        \max_{1\le i\le n}|\eps_{n,i}|\longrightarrow0,
\end{equation}
where $a_n>0$ and $f$ is continuous and bounded above and below by positive
constants.  Then, for every $1\le p<\infty$,
\begin{equation}\label{eq:main-shape-conv}
 \norm{
    \sup_{(x,y)\in[0,1]^2}
    \abs{\frac{\Lambda_n(x,y)}{\sqrt n}
          -\calU_{\rho_f}(x,y)}
 }_{L^p}
 \longrightarrow0.
\end{equation}
For every $\delta>0$ there are $c_\delta,C_\delta>0$ such that
\begin{equation}\label{eq:main-exp-conc}
 \Pp\left(
    \sup_{(x,y)\in[0,1]^2}
    \abs{\frac{\Lambda_n(x,y)}{\sqrt n}
          -\calU_{\rho_f}(x,y)}>\delta
 \right)
 \le C_\delta e^{-c_\delta\sqrt n}
\end{equation}
for all sufficiently large $n$.  Consequently the convergence is almost sure
under every coupling of the arrays across $n$.

In particular,
\begin{equation}\label{eq:main-lis-conv}
        \frac{\LIS(\sigma_n)}{\sqrt n}
        \longrightarrow
        \calL_+(f):=\calU_{\rho_f}(1,1)
\end{equation}
in every finite $L^p$ and almost surely, and
\[
        \frac{\Ee\LIS(\sigma_n)}{\sqrt n}
        \longrightarrow\calL_+(f).
\]
\end{theorem}

The theorem remains valid when the profile approximation fails on a
sublinear set of rows.

\begin{theorem}[Robust profile theorem]\label{thm:robust-luce}
Let $f\in C[0,1]$ be bounded above and below by positive constants.  After
removing a common rate factor, suppose
\begin{equation}\label{eq:comparable-rates}
        0<\underline\theta\le\theta_{n,i}
        \le\overline\theta<\infty
\end{equation}
uniformly in $n,i$, and suppose that for every $\eta>0$,
\begin{equation}\label{eq:profile-density-conv}
 \frac1n\#\left\{1\le i\le n:
       |\theta_{n,i}-f(i/n)|>\eta\right\}
 \longrightarrow0.
\end{equation}
Then all conclusions of Theorem~\ref{thm:main-luce} remain valid with the same
limit $\calU_{\rho_f}$.
\end{theorem}

\subsection{Endpoint profiles and the Sukhatme law}

For a continuous profile $f$ that is positive on $[0,1)$ and vanishes at
$1$, the same formula \eqref{eq:rho-def-intro} defines a possibly singular
permuton density on $(0,1)^2$.  Its terminal Hammersley constant is defined by
\begin{equation}\label{eq:singular-functional-intro}
        \calL_+^{\mathrm{sing}}(f)
        =\lim_{\delta\downarrow0}
          \calU_{\rho_f}(1-\delta,1).
\end{equation}
The endpoint theorem below gives transparent finite-array hypotheses under
which this limit is finite and is the LIS constant.

\begin{corollary}[Sukhatme weights]\label{cor:intro-sukhatme}
Let $\sigma_n$ be Luce distributed with
\[
        \theta_{n,i}=n-i+1,
        \qquad 1\le i\le n.
\]
Then, for every $1\le p<\infty$,
\begin{equation}\label{eq:sukhatme-intro}
        \frac{\LIS(\sigma_n)}{\sqrt n}
        \longrightarrow\calL_{\mathrm{Suk}}
\end{equation}
in $L^p$, with fixed-deviation probability $O(e^{-c\sqrt n})$ and almost
surely under arbitrary couplings.  Here
\[
        \calL_{\mathrm{Suk}}
        =\calL_+^{\mathrm{sing}}(f),
        \qquad f(x)=1-x,
\]
and the corresponding density is
\begin{equation}\label{eq:suk-rho-intro}
 \rho_{\mathrm{Suk}}(x,y)
 =\frac{(1-x)\tau(y)^2e^{x\tau(y)}}
        {e^{\tau(y)}-\tau(y)-1},
\end{equation}
where $\tau$ is the inverse of
$A(t)=1-(1-e^{-t})/t$.
\end{corollary}

Corollary~\ref{cor:sukhatme-orientations} also gives the decreasing-subsequence
law and the corresponding laws for the reversed rates $\theta_{n,i}=i$.
The two constants are exchanged under label reversal.
Proposition~\ref{prop:sukhatme-truncation} gives a multiplicative truncation
bound for both constants, and Proposition~\ref{prop:sukhatme-explicit-lower}
proves $\calL_{\mathrm{Suk}}>3.0602$ by an explicit trial path.

\subsection{Additional consequences}

The row-stratified theorem proves more than convergence of lengths.  Every
near-maximal chain is exponentially likely to lie close to the compact set of
maximizing curves; see Theorem~\ref{thm:localization}.  Vertical reflection
gives a longest-decreasing-subsequence law.  Let $\calA_-(1,1)$ be the class
of absolutely continuous paths from $(0,1)$ to $(1,0)$ with first coordinate
nondecreasing and second coordinate nonincreasing, and define
\begin{equation}\label{eq:Lminus-def}
 \calL_-(f)=2\sup_{\gamma\in\calA_-(1,1)}
 \int_0^1\sqrt{\rho_f(\gamma_1(t),\gamma_2(t))
                 \gamma_1'(t)(-\gamma_2'(t))}\,\dd t,
\end{equation}
with endpoint actions defined by monotone truncation.  If
$\kappa_f=M_f/m_f$, then
\begin{equation}\label{eq:universal-bounds-intro}
        \frac2{\sqrt{\kappa_f}}
        \le \calL_+(f),\calL_-(f)
        \le2\sqrt{\kappa_f},
        \qquad
        \calL_+(f)\calL_-(f)\ge1.
\end{equation}
The constants are locally Lipschitz in the positive profile.  These
consequences are proved in Section~\ref{sec:consequences}.

\section{Row-stratified point clouds}

\subsection{The model and its variational functional}

Let $\partialmark$ be a cemetery state.  For each $n$, let
\[
        Y_{n,1},\dots,Y_{n,n}
\]
be independent random variables taking values in $[0,1]\cup\{\partialmark\}$.
Assume that the restriction of the law of $Y_{n,i}$ to $[0,1]$ has a
subdensity $g_{n,i}$; thus
\[
        g_{n,i}\ge0,
        \qquad
        \int_0^1g_{n,i}(y)\,\dd y\le1,
\]
and the remaining mass is placed at $\partialmark$.  The associated point
cloud contains $(i/n,Y_{n,i})$ when $Y_{n,i}\ne\partialmark$.
We call the deterministic vertical fibre $\{i/n\}\times[0,1]$ row $i$.

For $x,y\in[0,1]$, define $L_n(x,y)$ to be the largest $k$ for which
\[
        1\le i_1<\cdots<i_k\le\lfloor nx\rfloor,
        \qquad
        0\le Y_{n,i_1}<\cdots<Y_{n,i_k}\le y.
\]
Ties have probability zero under the density assumption.

Let $q\in C([0,1]^2)$ be nonnegative with
\begin{equation}\label{eq:q-row-mass}
        \int_0^1q(x,y)\,\dd y\le1
        \qquad (0\le x\le1).
\end{equation}
For $(x,y)\in[0,1]^2$, define
\begin{equation}\label{eq:Uq-def}
 \calU_q(x,y)=
 2\sup_{\gamma\in\calA_+(x,y)}
 \int_0^1
 \sqrt{q(\gamma_1(t),\gamma_2(t))
       \gamma_1'(t)\gamma_2'(t)}\,\dd t.
\end{equation}
The value is unchanged by absolutely continuous increasing
reparameterizations of the path.

\begin{lemma}[Basic properties of the value function]\label{lem:U-basic}
Let $q,\widetilde q$ be bounded nonnegative functions on the square.
Then
\begin{equation}\label{eq:U-bound}
        0\le\calU_q(x,y)\le2\sqrt{\norm q_\infty xy}.
\end{equation}
Moreover,
\begin{equation}\label{eq:U-sqrt-stability}
 \sup_{(x,y)\in[0,1]^2}
 |\calU_q(x,y)-\calU_{\widetilde q}(x,y)|
 \le2\norm{\sqrt q-\sqrt{\widetilde q}}_\infty.
\end{equation}
If $q\le M$, then for $x\le x'$ and $y\le y'$,
\begin{equation}\label{eq:U-modulus}
 0\le\calU_q(x',y')-\calU_q(x,y)
 \le2\sqrt M\left(
       \sqrt{(x'-x)y'}+\sqrt{x'(y'-y)}\right).
\end{equation}
In particular, $\calU_q$ is uniformly continuous.
\end{lemma}

\begin{proof}
For every admissible path,
\[
 \int_0^1\sqrt{\gamma_1'(t)\gamma_2'(t)}\,\dd t
 \le\sqrt{\int_0^1\gamma_1'(t)\,\dd t}
      \sqrt{\int_0^1\gamma_2'(t)\,\dd t}
 =\sqrt{xy}.
\]
This proves \eqref{eq:U-bound}.  Comparing the two actions along the same path
and using the same inequality proves \eqref{eq:U-sqrt-stability}.

To prove \eqref{eq:U-modulus}, take a path to $(x',y')$.  Before its first
exit from $[0,x]\times[0,y]$, its action is at most $\calU_q(x,y)$, since the
initial segment can be completed to $(x,y)$ by zero-action horizontal and
vertical pieces.  If the path exits through the right side, the remaining
horizontal increment is at most $x'-x$ and its vertical increment is at most
$y'$.  If it exits through the top, the corresponding increments are at most
$x'$ and $y'-y$.  Applying \eqref{eq:U-bound} to the remaining segment and
using the sum of the two possible bounds gives \eqref{eq:U-modulus}.
\end{proof}

\subsection{The stratified shape theorem}

\begin{theorem}[Row-stratified Hammersley shape theorem]
\label{thm:stratified-shape}
Assume
\begin{equation}\label{eq:uniform-density-conv}
        \eta_n:=
        \max_{1\le i\le n}
        \norm{g_{n,i}-q(i/n,\cdot)}_\infty
        \longrightarrow0,
\end{equation}
where $q\in C([0,1]^2)$ is nonnegative and satisfies
\eqref{eq:q-row-mass}.  Then, for every $1\le p<\infty$,
\begin{equation}\label{eq:stratified-Lp}
 \norm{
 \sup_{(x,y)\in[0,1]^2}
 \abs{\frac{L_n(x,y)}{\sqrt n}-\calU_q(x,y)}
 }_{L^p}
 \longrightarrow0.
\end{equation}
For every $\delta>0$ there are $c_\delta,C_\delta>0$ such that
\begin{equation}\label{eq:stratified-exp}
 \Pp\left(
 \sup_{(x,y)\in[0,1]^2}
 \abs{\frac{L_n(x,y)}{\sqrt n}-\calU_q(x,y)}>\delta
 \right)
 \le C_\delta e^{-c_\delta\sqrt n}
\end{equation}
for all sufficiently large $n$.  Consequently the convergence is almost sure
under every coupling of the arrays across $n$.
\end{theorem}

The proof occupies Sections~\ref{sec:homogeneous-input}--\ref{sec:shape-proof}.
The probabilistic inputs beyond elementary concentration are the classical
uniform-permutation LIS law and Talagrand's median inequality.

\section{Homogeneous input and a local decomposition}
\label{sec:homogeneous-input}

\subsection{Uniform permutations}

Let $H_N$ denote the LIS length of a uniformly random permutation of
$\{1,\dots,N\}$, with $H_0=0$.
For $N$ independent uniform points in a nondegenerate rectangle, sorting by
the first coordinate leaves independent uniform second coordinates.  Their
relative ranks form a uniform permutation, so the longest-chain length has
law $H_N$.  The same rank argument applies to independent identically
distributed continuous heights on any $N$ ordered deterministic fibres.

\begin{lemma}[Uniform-permutation input]\label{lem:uniform-lis}
For every $\delta>0$ there are constants $c_\delta,C_\delta>0$ such that
\begin{equation}\label{eq:uniform-lis-tail}
 \Pp\left(|H_N-2\sqrt N|>\delta\sqrt N\right)
 \le C_\delta e^{-c_\delta\sqrt N}
 \qquad(N\ge1).
\end{equation}
In particular, $H_N/\sqrt N\to2$ in probability.
\end{lemma}

\begin{proof}
Represent $H_N$ as the LIS length of $N$ independent uniform heights.
The Logan--Shepp and Vershik--Kerov theorem
\cite{LoganShepp1977,VershikKerov1977} gives $H_N/\sqrt N\to2$ in
probability.  Consequently, any median $m_N$ satisfies $m_N/\sqrt N\to2$.
Talagrand's median inequalities \cite[Theorem~7.1.2]{Talagrand1995} give,
for every $u>0$,
\begin{align*}
 \Pp(H_N\ge m_N+u)
 &\le2\exp\left\{-\frac{u^2}{4(m_N+u)}\right\},\\
 \Pp(H_N\le m_N-u)
 &\le2\exp\left\{-\frac{u^2}{4m_N}\right\}.
\end{align*}
For sufficiently large $N$,
$|m_N-2\sqrt N|\le(\delta/2)\sqrt N$ and $m_N\le3\sqrt N$.
Taking $u=(\delta/2)\sqrt N$ proves
\eqref{eq:uniform-lis-tail}; enlarging $C_\delta$ covers the finitely many
remaining values of $N$.
\end{proof}

Sharper nontrivial lower- and upper-tail large-deviation asymptotics have
speeds $N$ and $\sqrt N$, respectively
\cite{DeuschelZeitouni1999}.  The common bound
\eqref{eq:uniform-lis-tail} suffices for the arguments below.

\begin{lemma}[A first-moment chain bound]\label{lem:first-moment-chain}
Suppose a point cloud is generated by independent row variables, with at most
one point in each of $r$ ordered rows, and each row subdensity on an interval
$J$ is bounded by $D$.  If $R$ is the
longest chain using these rows and heights in $J$, then
\begin{equation}\label{eq:first-moment-chain}
        \Pp(R\ge k)
        \le \binom rk\frac{(D|J|)^k}{k!}
        \le\frac{(rD|J|)^k}{(k!)^2}.
\end{equation}
\end{lemma}

\begin{proof}
For fixed increasing row indices, the joint density is at most $D^k$.
The region of strictly increasing heights in $J^k$ has volume $|J|^k/k!$.
Summing over the $\binom rk$ possible row sets proves the claim.
\end{proof}

\subsection{Common core and residual proposal cloud}

Fix an integer $m\ge1$ and put
\[
        I_j=\left(\frac{j-1}{m},\frac jm\right],
        \qquad 1\le j\le m,
        \qquad \alpha=\frac1m.
\]
Define
\begin{equation}\label{eq:h-omega-def}
 h_j(y)=\min_{x\in\overline I_j}q(x,y),
 \qquad
 \omega_m=
 \sup_{\substack{x,x'\in[0,1],\ |x-x'|\le1/m\\y\in[0,1]}}
 |q(x,y)-q(x',y)|.
\end{equation}
Then $\omega_m\to0$ by uniform continuity.  Put
\begin{equation}\label{eq:hjn-def}
        h_{j,n}(y)=(h_j(y)-\eta_n)_+.
\end{equation}
If $i/n\in I_j$, then
\begin{equation}\label{eq:core-residual-density-bounds}
        0\le h_{j,n}(y)\le g_{n,i}(y),
        \qquad
        0\le g_{n,i}(y)-h_{j,n}(y)
        \le \omega_m+2\eta_n.
\end{equation}

For a vertical interval $J\subset[0,1]$, let $L_{n,j}(J)$ be the longest chain
formed by points in $I_j\times J$.

\begin{proposition}[Local core--residual estimate]
\label{prop:local-core-residual}
Fix $m$ so large that $\omega_m<1/2$.  For fixed $j,J$, put
\[
        H_j(J)=\int_Jh_j(y)\,\dd y.
\]
Then, for every $\delta>0$, with probability at least
$1-Ce^{-c\sqrt n}$ for all sufficiently large $n$,
\begin{align}
 L_{n,j}(J)
 &\ge 2\sqrt{\alpha nH_j(J)}-\delta\sqrt n,
 \label{eq:local-lower}\\
 L_{n,j}(J)
 &\le 2\sqrt{\alpha nH_j(J)}
       +2\sqrt{\alpha n\omega_m|J|}
       +\delta\sqrt n.
 \label{eq:local-upper}
\end{align}
The constants may depend on $m,j,J,\delta,q$ but not on $n$.
\end{proposition}

\begin{proof}
For every row $i/n\in I_j$, mark a point as belonging to the common core in
$J$ by prescribing the joint subdensity
\[
        \Pp(C_i=1,\ Y_{n,i}\in\dd y)
        =h_{j,n}(y)\one_J(y)\,\dd y.
\]
This is legitimate by \eqref{eq:core-residual-density-bounds}, and the marking
can be performed independently row by row.  The core indicators are i.i.d.
Bernoulli with parameter
\[
        p_{j,n}(J)=\int_Jh_{j,n}(y)\,\dd y.
\]
Conditional on being in the core, all heights have the same continuous
density $h_{j,n}\one_J/p_{j,n}(J)$.  Conditional on the number $K$ of core
points, their relative vertical order is therefore an exact uniform random
permutation of size $K$.  Since the number of rows in $I_j$ is
$\alpha n+O(1)$ and $p_{j,n}(J)\to H_j(J)$, Lemma~\ref{lem:uniform-lis},
combined with a binomial Chernoff bound, gives both upper and lower estimates
\[
        L_{n,j}^{\mathrm{core}}(J)
        =2\sqrt{\alpha nH_j(J)}+o(\sqrt n)
\]
with failure probability $O(e^{-c\sqrt n})$.  If $H_j(J)=0$, continuity gives
$h_j=0$ on $J$, so the core is empty and the same statement holds.

It remains to control residual points in $J$.  By
\eqref{eq:core-residual-density-bounds}, their row subdensities are bounded by
\[
        d_{m,n}=\omega_m+2\eta_n.
\]
Since $\omega_m<1/2$ and $\eta_n\to0$, one has
$d_{m,n}|J|\le1$ for all sufficiently large $n$.  The residual points may
therefore be generated by thinning an independent proposal cloud that, in
each row, places a point uniformly in $J$ with probability $d_{m,n}|J|$.
Independence is required across rows; the proposal cloud need not be
independent of the common core.
Conditional on the proposal count, the proposal heights have a uniform random
relative order.  Lemma~\ref{lem:uniform-lis} and a
binomial Chernoff bound therefore give
\[
        L_{n,j}^{\mathrm{res}}(J)
        \le2\sqrt{\alpha n\omega_m|J|}+o(\sqrt n)
\]
with failure probability $O(e^{-c\sqrt n})$.  When $\omega_m=0$, use
Lemma~\ref{lem:first-moment-chain} and $d_{m,n}\to0$ instead.

Every chain splits into its core and residual subchains, so
\[
        L_{n,j}(J)
        \le L_{n,j}^{\mathrm{core}}(J)
           +L_{n,j}^{\mathrm{res}}(J),
\]
whereas the core itself supplies a lower bound.  This proves
\eqref{eq:local-lower}--\eqref{eq:local-upper}.
\end{proof}



\section{Proof of the stratified shape theorem}
\label{sec:shape-proof}

\subsection{The lower bound}

For a nondecreasing partition
\[
        0=z_0\le z_1\le\cdots\le z_m=1,
\]
define
\begin{equation}\label{eq:Fm-def}
 F_m=2\sup_{0=z_0\le\cdots\le z_m=1}
 \sum_{j=1}^m
 \sqrt{\frac1m\int_{z_{j-1}}^{z_j}h_j(y)\,\dd y}.
\end{equation}
Inside column $I_j$, choose a longest common-core chain in
$I_j\times(z_{j-1},z_j)$.  These $m$ chains concatenate because both their
horizontal and vertical ranges are ordered.  Proposition
\ref{prop:local-core-residual}, applied to the finitely many rectangles, gives
\begin{equation}\label{eq:global-liminf}
        \liminf_{n\to\infty}\frac{L_n(1,1)}{\sqrt n}
        \ge F_m
\end{equation}
in probability, with fixed-deviation failure probability $O(e^{-c\sqrt n})$.

\subsection{The upper bound and random exit heights}

Take a longest chain and let $s_j$ be the largest height attained by this
chain in the first $j$ columns, using $s_j=s_{j-1}$ when column $j$ contributes
no point, and setting $s_0=0$, $s_m=1$.  Every chain point in column $j$ lies
in
\[
        I_j\times[s_{j-1},s_j].
\]
The intervals are random, so we discretize them.  Fix $r\ge1$ and set
\[
 J_j^{(r)}=
 \left[\frac{\lfloor rs_{j-1}\rfloor}{r},
       \frac{\lceil rs_j\rceil}{r}\right].
\]
For fixed $m,r$, only finitely many such intervals occur.  Proposition
\ref{prop:local-core-residual} therefore holds simultaneously for all of them.
Writing $M=\norm q_\infty$ and
$\widehat s_j=\lceil rs_j\rceil/r$, the intervals
$[\widehat s_{j-1},\widehat s_j]$ form a genuine partition.  The lower
overhang of $J_j^{(r)}$ has length at most $1/r$, and hence
\begin{equation}\label{eq:core-mesh-bound}
 \sum_{j=1}^m
 2\sqrt{\frac1m\int_{J_j^{(r)}}h_j}
 \le F_m+2\sqrt{\frac{mM}{r}}.
\end{equation}
Also
\[
        \sum_{j=1}^m|J_j^{(r)}|
        \le1+\frac{2m}{r}.
\]
Thus Cauchy--Schwarz yields
\begin{equation}\label{eq:residual-mesh-bound}
 \sum_{j=1}^m
 2\sqrt{\frac{\omega_m}{m}|J_j^{(r)}|}
 \le2\sqrt{\omega_m\left(1+\frac{2m}{r}\right)}.
\end{equation}
Summing the local upper bounds, then letting $r\to\infty$, gives
\begin{equation}\label{eq:global-limsup}
        \limsup_{n\to\infty}\frac{L_n(1,1)}{\sqrt n}
        \le F_m+2\sqrt{\omega_m}
\end{equation}
in probability, again with fixed-deviation failure probability
$O(e^{-c\sqrt n})$ for each fixed $m$.

\subsection{Identification of the discrete variational problem}

Define a piecewise constant density
\begin{equation}\label{eq:qm-def}
        q_m(x,y)=h_j(y),
        \qquad x\in I_j.
\end{equation}
Then
\begin{equation}\label{eq:q-qm-bound}
        \norm{q-q_m}_\infty\le\omega_m.
\end{equation}

\begin{lemma}[Exact block variational formula]\label{lem:block-variational}
For every $m$,
\begin{equation}\label{eq:Fm-equals-Uqm}
        F_m=\calU_{q_m}(1,1).
\end{equation}
Consequently,
\begin{equation}\label{eq:Fm-to-Uq}
        F_m\longrightarrow\calU_q(1,1).
\end{equation}
\end{lemma}

\begin{proof}
Consider an admissible path and its passage through column $I_j$.  Suppose its
vertical coordinate rises there from $a$ to $b$.  Cauchy--Schwarz gives
\begin{equation}\label{eq:column-CS}
 \int_{\gamma\cap(I_j\times[0,1])}
 \sqrt{h_j(y)x'y'}\,\dd t
 \le
 \sqrt{\frac1m\int_a^bh_j(y)\,\dd y}.
\end{equation}
Summing over columns proves $\calU_{q_m}(1,1)\le F_m$.

Conversely, fix a partition $z_0,\dots,z_m$ and put
\[
        p_j=\int_{z_{j-1}}^{z_j}h_j(y)\,\dd y.
\]
When $p_j>0$, traverse column $j$ along the absolutely continuous curve
specified by
\[
 x(y)=\frac{j-1}{m}
      +\frac1m\frac{\int_{z_{j-1}}^yh_j(s)\,\dd s}{p_j},
 \qquad z_{j-1}\le y\le z_j.
\]
Equality holds in \eqref{eq:column-CS}, and the contribution to the action is
$2\sqrt{p_j/m}$.  If $p_j=0$, traverse the required horizontal and vertical
increments separately; their action is zero.  Concatenating the column paths
proves the reverse inequality in \eqref{eq:Fm-equals-Uqm}.

Finally, Lemma~\ref{lem:U-basic} and \eqref{eq:q-qm-bound} imply
\[
 |\calU_q(1,1)-\calU_{q_m}(1,1)|
 \le2\norm{\sqrt q-\sqrt{q_m}}_\infty
 \le2\sqrt{\omega_m},
\]
which proves \eqref{eq:Fm-to-Uq}.
\end{proof}

Combining \eqref{eq:global-liminf}, \eqref{eq:global-limsup},
\eqref{eq:Fm-to-Uq}, and $\omega_m\to0$ proves
\begin{equation}\label{eq:endpoint-prob-conv}
        \frac{L_n(1,1)}{\sqrt n}
        \longrightarrow\calU_q(1,1)
\end{equation}
in probability.  Choosing $m$ and then $r$ as functions only of a prescribed
fixed deviation $\delta$, and using the exponential estimates in
Proposition~\ref{prop:local-core-residual}, proves the endpoint form of
\eqref{eq:stratified-exp}.

\subsection{The full shape field}

Fix $(x,y)\in[0,1]^2$.  Apply the preceding argument to the rectangle
$[0,x]\times[0,y]$.  The last column may have width smaller than
$1/m$; treating it as one additional column changes every block value by at
most $O(m^{-1/2})$.  The lower construction concatenates core chains only
inside the rectangle.  In the upper construction, the random exit height in
each complete column is snapped to the same finite vertical mesh, while the
partial final column is bounded by Lemma~\ref{lem:first-moment-chain}.  Thus,
for every fixed $(x,y)$,
\begin{equation}\label{eq:pointwise-field}
        \frac{L_n(x,y)}{\sqrt n}
        \longrightarrow\calU_q(x,y)
\end{equation}
with fixed-deviation failure probability $O(e^{-c\sqrt n})$.

To pass to uniform convergence, choose a finite square grid of mesh $\zeta$.
If $(x_-,y_-)$ and $(x_+,y_+)$ are the lower-left and upper-right grid corners
surrounding $(x,y)$, monotonicity gives
\[
        L_n(x_-,y_- )\le L_n(x,y)\le L_n(x_+,y_+).
\]
Lemma~\ref{lem:U-basic} makes the corresponding oscillation of $\calU_q$
uniformly small when $\zeta$ is small.  A union bound over the fixed grid,
followed by the monotonicity sandwich, proves \eqref{eq:stratified-exp} and
uniform convergence in probability.

\subsection{Moments and almost sure convergence}

Let $M$ bound all row densities for sufficiently large $n$.  Applying
Lemma~\ref{lem:first-moment-chain} with $r=n$, $J=[0,1]$, and $D=M$ gives
\begin{equation}\label{eq:global-chain-tail}
        \Pp(L_n(1,1)\ge k)
        \le\frac{(nM)^k}{(k!)^2}.
\end{equation}
Stirling's formula shows that
$\{(L_n(1,1)/\sqrt n)^p:n\ge1\}$ is uniformly integrable for every finite
$p$.  Since the supremum norm error in Theorem~\ref{thm:stratified-shape} is
bounded by $L_n(1,1)/\sqrt n+2\sqrt M$, uniform convergence in probability
upgrades to \eqref{eq:stratified-Lp}.  The exponential estimate
\eqref{eq:stratified-exp} is summable in $n$, so Borel--Cantelli proves almost
sure convergence under every coupling.  This completes the proof of
Theorem~\ref{thm:stratified-shape}.


\section{Hamilton--Jacobi interpretation}

\begin{proposition}[Hamilton--Jacobi characterization]\label{prop:HJ}
For every nonnegative $q\in C([0,1]^2)$, the value function
\eqref{eq:Uq-def} is the unique continuous, coordinatewise nondecreasing
viscosity solution of
\begin{equation}\label{eq:HJ}
        (\calU_q)_x(\calU_q)_y=q
        \qquad\text{in }(0,1)^2,
\end{equation}
with incoming boundary values
\[
        \calU_q(x,0)=\calU_q(0,y)=0.
\]
No boundary data are prescribed on $\{x=1\}\cup\{y=1\}$.
\end{proposition}

\begin{proof}
Continuity, monotonicity, and the boundary values follow from
Lemma~\ref{lem:U-basic}.  The viscosity equation follows from
\cite[Theorem~2.5]{CalderEsedogluHero2014}, applied to the zero extension of
$q|_{(0,1)^2}$.  This extension does not change the action on the square:
segments on its boundary have zero action.  Their value function omits the
leading factor $2$ and solves $U_xU_y=q/4$; multiplying it by $2$ gives
\eqref{eq:HJ}.

For uniqueness, let $u,v$ be two solutions in the stated class.  Fix
$\eps>0$ and put
\[
 w_\eps(x,y)=v(x,y)+\frac{\eps}{1-x}+\frac{\eps}{1-y},
 \qquad (x,y)\in[0,1)^2.
\]
Every gradient $p$ of a smooth function touching $v$ from below at an
interior point has $p_1,p_2\ge0$, by coordinatewise monotonicity, and
$p_1p_2\ge q$, by the supersolution property.  Consequently every such
test for $w_\eps$ satisfies
\[
 \left(p_1+\frac{\eps}{(1-x)^2}\right)
 \left(p_2+\frac{\eps}{(1-y)^2}\right)
 \ge q(x,y)+\eps^2.
\]
Thus $w_\eps$ is a strict supersolution.  It exceeds $u$ on the incoming
axes and tends to infinity at the outgoing edges.  If
$\sup(u-w_\eps)>0$, maximize
$u(z)-w_\eps(\zeta)-\alpha|z-\zeta|^2/2$.
The barrier keeps $\zeta$ away from the outgoing edges, and
$|z-\zeta|\to0$ as $\alpha\to\infty$.  The strict incoming boundary
inequality then keeps both maximizing points in the interior for large
$\alpha$.  The two viscosity inequalities, with their common test
gradient, imply
\[
        \eps^2\le q(z)-q(\zeta)\longrightarrow0,
\]
a contradiction.  Hence $u\le w_\eps$.  Letting $\eps\downarrow0$,
using continuity at the outgoing edges, and interchanging $u,v$ proves
uniqueness.  This argument requires no strict positivity of $q$.
\end{proof}

Theorem~\ref{thm:stratified-shape} therefore supplies a row-stratified
probabilistic approximation to this Hamilton--Jacobi field, with exactly
one opportunity for a point in each deterministic row.


\section{Localization of near-maximal chains}
\label{sec:localization}

The finite-block upper bound also locates extremal chains.  We work on the
compact space $\Gamma$ of continuous paths
$\gamma=(\gamma_1,\gamma_2):[0,2]\to[0,1]^2$ satisfying
\[
 \gamma(0)=(0,0),\qquad \gamma(2)=(1,1),\qquad
 \gamma_1,\gamma_2\text{ are nondecreasing},\qquad
 \gamma_1(s)+\gamma_2(s)=s.
\]
Every continuous monotone path from $(0,0)$ to $(1,1)$ has this canonical
parameterization after deleting pauses.  Each coordinate is $1$-Lipschitz.
Equip $\Gamma$ with the metric
\[
 d_\Gamma(\gamma,\widetilde\gamma)
 =\sup_{0\le s\le2}\max_{j=1,2}
       |\gamma_j(s)-\widetilde\gamma_j(s)|.
\]
Arzel\`a--Ascoli shows that $\Gamma$ is compact.  Define
\begin{equation}\label{eq:path-action}
 I_q(\gamma)=2\int_0^2
 \sqrt{q(\gamma(s))\gamma_1'(s)\gamma_2'(s)}\,\dd s
\end{equation}
and let
\[
 \mathscr M_q
 =\{\gamma\in\Gamma:I_q(\gamma)=\calU_q(1,1)\}.
\]
The supremum over $\Gamma$ equals $\calU_q(1,1)$ by invariance of the action
under monotone reparameterization.  Distances to subsets of $\Gamma$ are
taken in $d_\Gamma$.

\begin{lemma}[Compactness and an action gap]\label{lem:action-gap}
If $q$ is continuous and nonnegative, then $I_q$ is upper semicontinuous on
$\Gamma$.  Consequently $\mathscr M_q$ is nonempty and compact.  For every
$\eps>0$ there is $\eta_\eps>0$ such that
\begin{equation}\label{eq:action-gap}
 \sup\{I_q(\gamma):\dist(\gamma,\mathscr M_q)\ge\eps\}
 \le\calU_q(1,1)-4\eta_\eps,
\end{equation}
where the assertion is vacuous if the set on the left is empty.
\end{lemma}

\begin{proof}
For a partition $\Pi:0=s_0<\cdots<s_r=2$, let $R_\ell(\gamma)$ be the
coordinate rectangle spanned by $\gamma(s_{\ell-1})$ and $\gamma(s_\ell)$,
put $M_\ell(\gamma)=\max_{R_\ell(\gamma)}q$, and define
\[
 S_\Pi(\gamma)=2\sum_{\ell=1}^r
 \sqrt{M_\ell(\gamma)
       \Delta_\ell\gamma_1\Delta_\ell\gamma_2}.
\]
Cauchy--Schwarz gives $I_q(\gamma)\le S_\Pi(\gamma)$, and refinement can
only decrease these upper sums.  Along a nested refining sequence with mesh
tending to zero, uniform continuity of $q$ and Lebesgue differentiation of
$\gamma_1',\gamma_2'$ give $S_\Pi(\gamma)\downarrow I_q(\gamma)$.  Hence
\[
 I_q(\gamma)=\inf_\Pi S_\Pi(\gamma).
\]
For fixed $\Pi$, the map $S_\Pi$ is continuous, since the maximum of a
continuous function over a coordinate rectangle depends continuously on its
endpoints.  Thus $I_q$ is upper semicontinuous.  Compactness gives existence
of maximizers and the strict gap away from their maximizing set.
\end{proof}

For a finite increasing chain $\mathcal C$, let
\[
 \Gamma(\mathcal C)
 =\{\gamma\in\Gamma:\mathcal C\subseteq\gamma([0,2])\}
\]
be its set of monotone interpolants.  This set contains the canonical
polygonal path $\gamma_{\mathcal C}$ obtained by adjoining $(0,0),(1,1)$
and joining successive vertices by straight segments.  For a nonempty
closed set $F\subseteq\Gamma$, define
\[
 M_n(F)=\max\{|\mathcal C|:
       \mathcal C\text{ is a chain in the array and }
       \Gamma(\mathcal C)\cap F\ne\varnothing\},
\]
with value $0$ if there is no such chain.  This is a measurable random
variable: for each of the finitely many choices of rows, the existence of an
interpolant in compact $F$ is a closed condition on the corresponding
heights.

\begin{lemma}[Restricted upper shape bound]\label{lem:restricted-upper}
Under the hypotheses of Theorem~\ref{thm:stratified-shape}, for every
nonempty closed $F\subseteq\Gamma$ and every $\delta>0$ there are $c,C>0$
such that
\begin{equation}\label{eq:restricted-upper}
 \Pp\left(
 M_n(F)>\left(\sup_{\gamma\in F}I_q(\gamma)+\delta\right)\sqrt n
 \right)
 \le Ce^{-c\sqrt n}
\end{equation}
for all sufficiently large $n$.
\end{lemma}

\begin{proof}
Put $a=\sup_F I_q$ and $M=\norm q_\infty$.  For each $\gamma\in F$, choose
a partition $\Pi_\gamma$ with
\[
 S_{\Pi_\gamma}(\gamma)<a+\delta/8.
\]
Refine it so that every interval has length less than $(4(M+1))^{-1}$.
Continuity gives a neighborhood $V_\gamma$ on which
$S_{\Pi_\gamma}<a+\delta/4$.  Compactness of $F$ gives a finite subcover.

Fix one of these neighborhoods and its partition
$\Pi:0=s_0<\cdots<s_r=2$.  Suppose $\mathcal C$ has an interpolant
$\gamma$ in that neighborhood.  Assign each chain vertex $(x,y)$ to the
interval containing its canonical parameter $s=x+y$, using
$[s_{\ell-1},s_\ell)$ and making the last interval closed on the right.
Every assigned vertex lies in the coordinate rectangle spanned by
$\gamma(s_{\ell-1})$ and $\gamma(s_\ell)$.

Snap each lower-left corner down and each upper-right corner up to a fixed
square grid of sufficiently small mesh $\zeta$.  There are finitely many
resulting rectangle sequences.  The mesh can be chosen so that every
enlarged vertical side has length at most $(2(M+1))^{-1}$ and, uniformly
over sequences arising from this neighborhood,
\begin{equation}\label{eq:snapped-upper-sum}
 2\sum_\ell
 \sqrt{\bigl(\sup_{R_\ell}q\bigr)
        |\operatorname{proj}_xR_\ell|
        |\operatorname{proj}_yR_\ell|}
 \le a+\delta/2.
\end{equation}
This follows from uniform continuity of $q$ and continuity of the upper sums
in their finitely many rectangle endpoints.

For one fixed snapped rectangle
$R=[u_0,u_1]\times[v_0,v_1]$, put
\[
 D_n(R)=\sup_Rq+\eta_n,
 \qquad
 r_n(R)=\#\{i:u_0\le i/n\le u_1\}
       =n(u_1-u_0)+O(1).
\]
Both vertical boundaries are included in this row count.  For large $n$,
$D_n(R)(v_1-v_0)\le1$, so each restricted row law is dominated by a proposal
which is present with probability $D_n(R)(v_1-v_0)$ and is uniform on that
interval when present.  Conditional on the proposal count, the relative
vertical order is a uniform permutation.  Lemma~\ref{lem:uniform-lis} and
a binomial Chernoff bound therefore imply, for every fixed $\xi>0$,
\begin{equation}\label{eq:rectangle-upper-tail}
 \Pp\left(
 L_n(R)>
 2\sqrt{nD_n(R)(u_1-u_0)(v_1-v_0)}+\xi\sqrt n
 \right)
 \le C_{R,\xi}e^{-c_{R,\xi}\sqrt n}.
\end{equation}
Here $L_n(R)$ counts chains in the closed rectangle.  If $\sup_Rq=0$, the
same estimate follows from Lemma~\ref{lem:first-moment-chain} and
$D_n(R)\to0$.  A zero-width rectangle contains at most one row.  In all
cases the $O(1)$ row-count discrepancy is absorbed by $\xi\sqrt n$.

Apply \eqref{eq:rectangle-upper-tail} to the finite collection of rectangles,
taking the sum of its $\xi$-errors smaller than $\delta/4$.  Each chain is
bounded by the sum of the chain lengths in its assigned rectangles.  Since
$D_n(R)\to\sup_Rq$ uniformly over this finite collection,
\eqref{eq:snapped-upper-sum} bounds that sum by $(a+\delta)\sqrt n$
outside an event of probability $Ce^{-c\sqrt n}$.  A final union bound over
the finite cover proves the result.  This argument uses only that the
interpolant contains the chain vertices.
\end{proof}

\begin{theorem}[Localization of near-maximal chains]
\label{thm:localization}
Under the hypotheses of Theorem~\ref{thm:stratified-shape}, for every
$\eps>0$ there are $\eta_\eps,c_\eps,C_\eps>0$ such that
\begin{equation}\label{eq:restricted-chain-tail}
 \Pp\left(
 \begin{array}{c}
 \text{there are a chain }\mathcal C
     \text{ and }\gamma\in\Gamma(\mathcal C)\text{ with}\cr
 |\mathcal C|\ge(\calU_q(1,1)-\eta_\eps)\sqrt n,
 \qquad\dist(\gamma,\mathscr M_q)\ge\eps
 \end{array}
 \right)
 \le C_\eps e^{-c_\eps\sqrt n}
\end{equation}
for all sufficiently large $n$.  In particular, with the same type of
probability bound, every monotone interpolant of every longest chain lies
in an $\eps$-neighborhood of $\mathscr M_q$.  If the maximizing path is
unique, these interpolants converge to it in probability and almost surely
under arbitrary couplings.
\end{theorem}

\begin{proof}
Let $F_\eps=\{\gamma\in\Gamma:
\dist(\gamma,\mathscr M_q)\ge\eps\}$.  If this set is empty, the assertions
are immediate.  Otherwise it is compact, and Lemma~\ref{lem:action-gap}
gives $\sup_{F_\eps}I_q\le\calU_q(1,1)-4\eta_\eps$ for some
$\eta_\eps>0$.  Lemma~\ref{lem:restricted-upper}, with
$\delta=\eta_\eps$, bounds every chain admitting an interpolant in
$F_\eps$ by $(\calU_q(1,1)-3\eta_\eps)\sqrt n$, outside the stated
exceptional event.  This proves \eqref{eq:restricted-chain-tail}.

The shape theorem supplies a chain of length at least
$(\calU_q(1,1)-2\eta_\eps)\sqrt n$ with the same type of failure
probability.  Thus no longest chain admits an interpolant in $F_\eps$
outside the union of these two events.  Borel--Cantelli, first for positive
rational $\eps$, gives the almost-sure assertion.
\end{proof}


\section{Rankings generated by independent scores}
\label{sec:independent-scores}

Let $X_{n,1},\dots,X_{n,n}$ be independent continuous random variables in
an interval $I\subseteq\Rr$, and let $\Psi:I\to(0,1)$ be a fixed
increasing homeomorphism.  Put $Z_{n,i}=\Psi(X_{n,i})$ and suppose that
$Z_{n,i}$ has density $g_{n,i}$.  Assume
\begin{equation}\label{eq:score-density-hyp}
 \max_i\norm{g_{n,i}-q(i/n,\cdot)}_\infty\longrightarrow0,
\end{equation}
where $q\in C([0,1]^2)$ is nonnegative and
$\int_0^1q(x,z)\,\dd z=1$ for every $x$.  Define
\begin{equation}\label{eq:mixture-cdf}
 b(z)=\int_0^1q(x,z)\,\dd x,
 \qquad Q(z)=\int_0^zb(s)\,\dd s.
\end{equation}
Assume $b(z)>0$ for $0<z<1$.  Then $Q$ is an increasing homeomorphism of
$[0,1]$, and the induced rank-coordinate density is
\begin{equation}\label{eq:score-permuton-density}
 \rho_q(x,y)
 =\frac{q(x,Q^{-1}(y))}{b(Q^{-1}(y))},
 \qquad 0<y<1.
\end{equation}
Both marginals of $\rho_q$ are uniform.  Its actions at the horizontal
edges are interpreted by monotone truncation from $0<y<1$.

\begin{lemma}[Vertical coordinate changes]\label{lem:vertical-path-map}
Let $\mathcal T_Q:\Gamma\to\Gamma$ map a path vertically by
$(x,z)\mapsto(x,Q(z))$ and then give its image the canonical
parameterization.  This map is a homeomorphism.  If
$\omega_Q$ is the modulus of continuity of $Q$, then
\begin{equation}\label{eq:vertical-path-continuity}
 d_\Gamma(\mathcal T_Q\gamma,\mathcal T_Q\widetilde\gamma)
 \le2\bigl(d+\omega_Q(d)\bigr),
 \qquad d=d_\Gamma(\gamma,\widetilde\gamma).
\end{equation}
Moreover,
\begin{equation}\label{eq:score-action-invariance}
 I_{\rho_q}(\mathcal T_Q\gamma)=I_q(\gamma),
 \qquad
 \calU_{\rho_q}(x,y)=\calU_q(x,Q^{-1}(y)).
\end{equation}
In particular all these actions are finite, and the rank-coordinate
maximizing set is the nonempty compact set
\begin{equation}\label{eq:score-maximizer-map}
 \mathscr M_{\rho_q}
 =\{\beta\in\Gamma:I_{\rho_q}(\beta)=\calU_{\rho_q}(1,1)\}
 =\mathcal T_Q(\mathscr M_q).
\end{equation}
\end{lemma}

\begin{proof}
For $\gamma\in\Gamma$, the function
$h_\gamma(s)=\gamma_1(s)+Q(\gamma_2(s))$ is an increasing homeomorphism
of $[0,2]$, and
\[
 (\mathcal T_Q\gamma)(h_\gamma(s))
   =(\gamma_1(s),Q(\gamma_2(s))).
\]
For two paths at distance $d$, their transformed points at the same old
parameter differ by at most $\max\{d,\omega_Q(d)\}$, and their new
parameters differ by at most $d+\omega_Q(d)$.  The canonical transformed
paths are $1$-Lipschitz, proving \eqref{eq:vertical-path-continuity}.
The inverse is $\mathcal T_{Q^{-1}}$, whose continuity follows by the same
argument.

The density identity is
\begin{equation}\label{eq:score-density-transform}
 \rho_q(x,Q(z))Q'(z)=q(x,z).
\end{equation}
On $\eps\le z\le1-\eps$, both $Q$ and $Q^{-1}$ are continuously
differentiable with bounded derivatives, so the chain rule and
reparameterization invariance identify the two path actions on that
portion.  The two omitted portions of the $q$-action are bounded by
$4\sqrt{\norm q_\infty\eps}$: each has vertical increment at most
$\eps$ and horizontal increment at most $1$.  Letting $\eps\downarrow0$
and using monotone truncation at both horizontal edges proves the first
identity in \eqref{eq:score-action-invariance}.  The same argument applies
to paths ending at any $(x,Q(z))$, and taking suprema gives its second
identity.  The homeomorphism and equality of actions then give
\eqref{eq:score-maximizer-map}.
\end{proof}

Let
\[
 \sigma_n(i)=1+\#\{j:X_{n,j}<X_{n,i}\}
\]
be the label-to-rank permutation, and define its restricted LIS field
$\Lambda_n$ as in the introduction.  For an increasing chain $\mathcal C$
in its normalized plot, write $\beta_{\mathcal C}$ for the canonical
polygonal path through its vertices and $(0,0),(1,1)$.

\begin{theorem}[Independent-score ranking theorem]
\label{thm:independent-score}
Under the hypotheses above, for every $1\le p<\infty$,
\begin{equation}\label{eq:score-shape-law}
 \norm{\sup_{x,y\in[0,1]}
 \abs{\frac{\Lambda_n(x,y)}{\sqrt n}-\calU_{\rho_q}(x,y)}}_{L^p}
 \longrightarrow0,
\end{equation}
with fixed-deviation probability $O(e^{-c\sqrt n})$ and almost-sure
convergence under arbitrary couplings.  In particular,
\[
 \frac{\LIS(\sigma_n)}{\sqrt n}
 \longrightarrow\calU_{\rho_q}(1,1)
\]
in every such $L^p$ and almost surely, and the normalized expectations have
the same limit.  For every $\eps>0$ there are
$\eta_\eps,c_\eps,C_\eps>0$ such that
\begin{equation}\label{eq:score-localization}
 \Pp\left(
 \begin{array}{c}
 \text{there exists a rank chain }\mathcal C\text{ with}\cr
 |\mathcal C|\ge(\calU_{\rho_q}(1,1)-\eta_\eps)\sqrt n,
 \quad\dist(\beta_{\mathcal C},\mathscr M_{\rho_q})\ge\eps
 \end{array}
 \right)
 \le C_\eps e^{-c_\eps\sqrt n}
\end{equation}
for all sufficiently large $n$.  Every longest-chain polygon localizes
around $\mathscr M_{\rho_q}$ with the same type of bound.  If the maximizing
path is unique, these polygons converge to it in probability and almost
surely under arbitrary couplings.
\end{theorem}

\begin{proof}
Theorem~\ref{thm:stratified-shape} applies to $(i/n,Z_{n,i})$.  Let
\[
 F_n(z)=\frac1n\sum_{i=1}^n\one_{\{Z_{n,i}\le z\}},
 \qquad
 \overline F_n(z)=\Ee F_n(z)
 =\frac1n\sum_{i=1}^n\int_0^zg_{n,i}(s)\,\dd s.
\]
Uniform density convergence and uniform Riemann-sum approximation give
\begin{equation}\label{eq:score-mean-cdf}
 \sup_{z\in[0,1]}|\overline F_n(z)-Q(z)|\longrightarrow0.
\end{equation}
For fixed $\delta>0$, choose a deterministic grid whose successive
$\overline F_n$-increments are at most $\delta/4$.  It has $O_\delta(1)$
points.  Hoeffding's inequality and monotonicity between grid points yield
\begin{equation}\label{eq:score-empirical-cdf}
 \Pp\left(\sup_z|F_n(z)-\overline F_n(z)|>\delta\right)
 \le C_\delta e^{-c_\delta n}.
\end{equation}
Ties have probability zero, so $\sigma_n(i)/n=F_n(Z_{n,i})$.  On
$\norm{F_n-Q}_\infty\le\delta$, the rank field is therefore sandwiched
between score-chain fields cut off at
$Q^{-1}((y-\delta-n^{-1})_+)$ and
$Q^{-1}((y+\delta)\wedge1)$.  Lemma~\ref{lem:vertical-path-map}, uniform
continuity of $\calU_q$ and $Q^{-1}$, and the stratified shape theorem give
uniform convergence and the asserted fixed-deviation bound.  Bounded row
densities give uniform integrability by
Lemma~\ref{lem:first-moment-chain}, so the convergence holds in every
$L^p$, $1\le p<\infty$.  Summability gives the almost-sure assertion.

To prove localization, fix $\eps>0$ and choose $0<\delta<\eps/4$.
For the score vertices of a rank chain $\mathcal C$, form the ideal rank
polygon $\beta^\circ_{\mathcal C}$ through
$(i/n,Q(Z_{n,i}))$, with endpoints $(0,0),(1,1)$.  On
$\norm{F_n-Q}_\infty\le\delta$,
\begin{equation}\label{eq:ideal-rank-polygon}
 d_\Gamma(\beta_{\mathcal C},\beta^\circ_{\mathcal C})\le2\delta
 \qquad\text{for every chain }\mathcal C.
\end{equation}
Indeed, linear interpolation at corresponding vertices leaves the first
coordinate unchanged and changes the second by at most $\delta$; the
canonical parameters consequently differ by at most $\delta$.  The
$1$-Lipschitz property of canonical paths gives the stated bound, including
after deletion of zero-length segments.

The path $\mathcal T_Q^{-1}\beta^\circ_{\mathcal C}$ is a monotone
interpolant through the exact score vertices.  If
$\dist(\beta_{\mathcal C},\mathscr M_{\rho_q})\ge\eps$, then
\eqref{eq:ideal-rank-polygon} places this interpolant in the compact set
\[
 F=\mathcal T_Q^{-1}
   \{\beta\in\Gamma:
          \dist(\beta,\mathscr M_{\rho_q})\ge\eps/2\}.
\]
This set is disjoint from $\mathscr M_q$.  If it is empty there is nothing
to prove.  Otherwise upper semicontinuity and compactness give
$\sup_F I_q\le\calU_q(1,1)-4\eta_\eps$ for some $\eta_\eps>0$.
Lemma~\ref{lem:restricted-upper}, with deviation $\eta_\eps$, bounds every
such chain by $(\calU_q(1,1)-3\eta_\eps)\sqrt n$ except with probability
$O(e^{-c\sqrt n})$.  Equations \eqref{eq:score-mean-cdf} and
\eqref{eq:score-empirical-cdf} bound the additional empirical-CDF exception
by $O(e^{-c'n})$ for large $n$.  This proves
\eqref{eq:score-localization}.  The lower shape bound and Borel--Cantelli
give the statements about longest chains and arbitrary couplings.
\end{proof}

\subsection{Gaussian Thurstone--Mosteller rankings}

The independent-score theorem includes the classical normal random-utility
model introduced by Thurstone and developed by Mosteller
\cite{Thurstone1927,Mosteller1951}.  Let
\begin{equation}\label{eq:thurstone-model}
        X_{n,i}=\mu_{n,i}+\sigma G_i,
        \qquad G_i\stackrel{\mathrm{i.i.d.}}\sim N(0,1),
\end{equation}
where $\sigma>0$ and
$\max_i|\mu_{n,i}-\mu(i/n)|\to0$ for a continuous bounded function $\mu$.
Choose $\tau>\sigma$ and let $H_\tau$ be the distribution function of
$N(0,\tau^2)$.  Taking $\Psi=H_\tau$, the limiting transformed density is
\begin{equation}\label{eq:thurstone-q}
 q_\mu(x,z)
 =\frac{\varphi_\sigma(H_\tau^{-1}(z)-\mu(x))}
        {\varphi_\tau(H_\tau^{-1}(z))},
\end{equation}
where $\varphi_s$ denotes the $N(0,s^2)$ density.

\begin{corollary}[Thurstone--Mosteller shape law]
\label{cor:thurstone}
For the ranking generated by \eqref{eq:thurstone-model}, the complete LIS
field satisfies Theorem~\ref{thm:independent-score} with $q=q_\mu$ and
permuton density given by \eqref{eq:score-permuton-density}.
\end{corollary}

\begin{proof}
Writing $t=H_\tau^{-1}(z)$, the logarithm of the ratio in
\eqref{eq:thurstone-q} has leading term
\[
 -\frac12\left(\frac1{\sigma^2}-\frac1{\tau^2}\right)t^2
 +\frac{\mu(x)}{\sigma^2}t+O(1).
\]
Since $\tau>\sigma$ and $\mu$ is bounded, $q_\mu(x,z)$ tends uniformly to
zero as $z\downarrow0$ or $z\uparrow1$.  It therefore extends to a bounded
continuous density on the closed square.  The derivative of the same density
ratio with respect to the location parameter is uniformly bounded when that
parameter ranges over a fixed compact set.  Hence
$\max_i|\mu_{n,i}-\mu(i/n)|\to0$ implies the uniform density convergence
required in \eqref{eq:score-density-hyp}, and
Theorem~\ref{thm:independent-score} applies.
\end{proof}


\section{Application to regular Luce profiles}
\label{sec:luce-application}

Remove the common factor $a_n$ in \eqref{eq:uniform-profile-hyp} and write
\[
        \lambda_{n,i}=f(i/n)+\eps_{n,i},
        \qquad
        T_{n,i}\sim\Exp(\lambda_{n,i})
\]
independently.  Since clock ties have probability zero,
\begin{equation}\label{eq:exact-lis-clock}
 \LIS(\sigma_n)
 =\max\left\{k:
    i_1<\cdots<i_k,\quad
    T_{n,i_1}<\cdots<T_{n,i_k}\right\}.
\end{equation}

Choose
\begin{equation}\label{eq:c-choice}
        0<c_0<\frac{m_f}{2}
\end{equation}
and compactify clock time by
\begin{equation}\label{eq:compactification}
        Z_{n,i}=1-e^{-c_0T_{n,i}}\in(0,1).
\end{equation}
The density of $Z_{n,i}$ is
\begin{equation}\label{eq:gin-luce}
        g_{n,i}(z)
        =\frac{\lambda_{n,i}}{c_0}
         (1-z)^{\lambda_{n,i}/c_0-1},
\end{equation}
and the limiting row density is
\begin{equation}\label{eq:qf-def}
        q_f(x,z)=\frac{f(x)}{c_0}
        (1-z)^{f(x)/c_0-1}.
\end{equation}
Because the exponents are uniformly bounded below by a positive constant,
$q_f$ is bounded and continuous on the closed square and
\begin{equation}\label{eq:luce-density-uniform}
        \max_i\norm{g_{n,i}-q_f(i/n,\cdot)}_\infty
        \longrightarrow0.
\end{equation}

The mixture distribution of the compactified score is
\begin{equation}\label{eq:luce-mixture-Q}
 Q_f(z)
 =A_f\left(-\frac1{c_0}\log(1-z)\right).
\end{equation}
Consequently
\[
 Q_f^{-1}(y)=1-e^{-c_0\tau_f(y)}
\]
and a direct substitution into \eqref{eq:score-permuton-density} gives
\[
 \rho_{q_f}(x,y)
 =\frac{f(x)e^{-f(x)\tau_f(y)}}{B_f(\tau_f(y))}
 =\rho_f(x,y).
\]
Theorem~\ref{thm:main-luce} is therefore the specialization of
Theorem~\ref{thm:independent-score} to exponential scores.

For later use, note the exact action identity
\begin{equation}\label{eq:luce-coordinate-invariance}
 \calU_{\rho_f}(x,y)
 =\calU_{q_f}\left(x,1-e^{-c_0\tau_f(y)}\right).
\end{equation}
The endpoint $y=1$ is obtained by monotone truncation.  In particular the
possibly unbounded Luce density has a finite Hammersley action whenever the
profile is bounded below by a positive constant.

\section{Robustness under sparse profile errors}

We prove Theorem~\ref{thm:robust-luce}.  Condition
\eqref{eq:profile-density-conv}, together with uniform boundedness, allows one
to choose a deterministic sequence $\delta_n\downarrow0$ such that
\begin{equation}\label{eq:exceptional-set-def}
 E_n=\{i:|\theta_{n,i}-f(i/n)|>\delta_n\}
 \qquad\text{satisfies}\qquad
 r_n:=|E_n|=o(n).
\end{equation}
Replace the rates on $E_n$ by $f(i/n)$ and leave all other rates unchanged.
The modified array converges uniformly to $f$ and hence satisfies
Theorem~\ref{thm:main-luce}.

Couple the original and modified clock arrays so that clocks outside $E_n$ are
identical.  If two real sequences agree outside a set $E$, then their LIS
lengths $L,L'$ satisfy
\begin{equation}\label{eq:sparse-replacement-lis}
        |L-L'|\le\max\{L(E),L'(E)\},
\end{equation}
where $L(E)$ and $L'(E)$ are the LIS lengths restricted to the exceptional
indices.  Indeed, delete the exceptional entries from a longest subsequence;
the remaining subsequence is present in both arrays, and the deleted entries
form an increasing subsequence on $E$.

For this section, decrease the compactification parameter $c_0$, if necessary,
so that
\[
        0<c_0<\frac12\min\{m_f,\underline\theta\}.
\]
All preceding coordinate identities remain valid for this choice.  After the
compactification \eqref{eq:compactification}, both the original and modified
exceptional-row densities are then bounded by a constant depending only on
$\underline\theta,\overline\theta,m_f,M_f$, and $c_0$.  Lemma
\ref{lem:first-moment-chain} therefore yields
\begin{equation}\label{eq:exceptional-chain-small}
        L(E_n),L'(E_n)=O_{\Pp}(\sqrt{r_n})
        =o_{\Pp}(\sqrt n),
\end{equation}
and, for every fixed $\delta>0$, with a bound
$O_\delta(e^{-c_\delta\sqrt n})$ for the event that either exceptional
chain exceeds $\delta\sqrt n$.

For the two-parameter permutation fields, changing the $r_n$ exceptional
clocks can shift the rank of a nonexceptional label by at most $r_n$.  After
deleting exceptional members of a chain, one obtains the deterministic
comparison
\begin{align}
 \Lambda_n(x,y)
 &\le \widetilde\Lambda_n
       \left(x,(y+r_n/n)\wedge1\right)+L(E_n),
 \label{eq:robust-field-upper}\\
 \widetilde\Lambda_n(x,y)
 &\le \Lambda_n
       \left(x,(y+r_n/n)\wedge1\right)+L'(E_n),
 \label{eq:robust-field-lower}
\end{align}
where tildes denote the modified array.  Uniform continuity of
$\calU_{\rho_f}$, together with $r_n/n\to0$ and
\eqref{eq:exceptional-chain-small}, transfers the full shape theorem,
moment convergence, concentration, and almost sure convergence from the
modified array to the original one.  This proves
Theorem~\ref{thm:robust-luce}.

\begin{remark}[Why permuton convergence alone is insufficient]
Theorem~\ref{thm:robust-luce} still uses two genuinely microscopic inputs:
uniform comparability of the rates and a sublinear exceptional set.  In an
arbitrary permutation model, one may plant an increasing block of size
$m_n$ with $\sqrt n\ll m_n\ll n$ without changing the permuton limit.  The
bounded-density estimate \eqref{eq:first-moment-chain} is what rules out this
mechanism for bounded-rate Luce rows.
\end{remark}

\section{Decreasing subsequences, bounds, and stability}\label{sec:consequences}

\subsection{Longest decreasing subsequences}

Replace every compactified height $Z_{n,i}$ by $1-Z_{n,i}$.  The transformed
row densities still satisfy the hypotheses of
Theorem~\ref{thm:stratified-shape}, and an increasing chain in the reflected
coordinates is a decreasing subsequence in the original coordinates.
In particular, if $q_f^\downarrow(x,z)=q_f(x,1-z)$, coordinate invariance gives
\begin{equation}\label{eq:Lminus-compactified}
        \calL_-(f)=\calU_{q_f^\downarrow}(1,1),
\end{equation}
independently of the admissible compactification parameter.  This proves the
following statement.

\begin{theorem}[Longest decreasing subsequences]\label{thm:lds}
Under either the hypotheses of Theorem~\ref{thm:main-luce} or those of
Theorem~\ref{thm:robust-luce},
\begin{equation}\label{eq:lds-conv}
        \frac{\LDS(\sigma_n)}{\sqrt n}
        \longrightarrow\calL_-(f)
\end{equation}
with fixed-deviation exponential concentration, almost sure convergence under
arbitrary couplings, and convergence in $L^p$ for every $1\le p<\infty$.
\end{theorem}

\subsection{Comparable-rate bounds}

\begin{proposition}[Universal bounds]\label{prop:universal-bounds}
Let $\kappa_f=M_f/m_f$.  Then
\begin{equation}\label{eq:universal-bounds}
        \frac2{\sqrt{\kappa_f}}
        \le \calL_+(f),\calL_-(f)
        \le2\sqrt{\kappa_f},
        \qquad
        \calL_+(f)\calL_-(f)\ge1.
\end{equation}
\end{proposition}

\begin{proof}
For each fixed $0<c_0<m_f$, uniform convergence of the rates eventually
keeps $\lambda_{n,i}>c_0$, so the continuous compactification and the coordinate
identities above remain valid.  The compactified density satisfies
\[
        q_f(x,z)\le\frac{M_f}{c_0}.
\]
Lemma~\ref{lem:U-basic} gives
\[
        \calL_+(f),\calL_-(f)
        \le2\sqrt{M_f/c_0}.
\]
Letting $c_0\uparrow m_f$ proves the upper bounds.  For the lower bounds, use
the row-independent subdensity
\[
 q_f(x,z)\ge h(z):=\frac{m_f}{c_0}(1-z)^{M_f/c_0-1},
 \qquad
 \int_0^1h(z)\,\dd z=\frac{m_f}{M_f}.
\]
For any nonnegative continuous $h$ depending only on the vertical coordinate,
Cauchy--Schwarz, with equality along
$x(z)=\int_0^z h/\int_0^1h$ when $\int h>0$, gives
\[
        \calU_h(1,1)=2\sqrt{\int_0^1h(z)\,\dd z}.
\]
Monotonicity of the action proves $\calL_+(f)\ge2/\sqrt{\kappa_f}$, and
vertical reflection gives the same bound for $\calL_-(f)$.
Every deterministic
permutation satisfies the Erd\H{o}s--Szekeres inequality
\[
        \LIS(\sigma)\LDS(\sigma)\ge n
\]
\cite{ErdosSzekeres1935}.  Passing to the limits in Theorems
\ref{thm:main-luce} and \ref{thm:lds} yields
$\calL_+(f)\calL_-(f)\ge1$.
\end{proof}

\begin{corollary}[Constant profile]\label{cor:constant-profile}
If $f$ is constant, then
\[
        \rho_f\equiv1,
        \qquad
        \calL_+(f)=\calL_-(f)=2.
\]
\end{corollary}

\subsection{Stability in the profile}

\begin{proposition}[Local Lipschitz stability]\label{prop:profile-stability}
Fix $0<m<M<\infty$.  There is $C_{m,M}<\infty$ such that whenever
$m\le f,g\le M$ are continuous,
\begin{equation}\label{eq:profile-stability}
        |\calL_+(f)-\calL_+(g)|
        +|\calL_-(f)-\calL_-(g)|
        \le C_{m,M}\norm{f-g}_\infty.
\end{equation}
\end{proposition}

\begin{proof}
Choose $c_0=m/2$.  For $a\in[m,M]$ and $z\in[0,1]$, put
\[
 s(a,z)=\sqrt{\frac{a}{c_0}}
        (1-z)^{(a/c_0-1)/2}.
\]
Because $(a/c_0-1)/2$ is bounded below by a positive constant, both
$s(a,z)$ and $s(a,z)|\log(1-z)|$ are uniformly bounded.  Differentiating with
respect to $a$ shows
\[
        \sup_{z\in[0,1]}
        |s(f(x),z)-s(g(x),z)|
        \le C_{m,M}\,|f(x)-g(x)|.
\]
Since $s(f(x),z)=\sqrt{q_f(x,z)}$, Lemma~\ref{lem:U-basic} gives the asserted
bound for $\calL_+$.  Vertical reflection gives the same estimate for
$\calL_-$.
\end{proof}



\section{Endpoint-degenerate Luce profiles}
\label{sec:endpoint}

We now allow the limiting rate profile to vanish at the terminal label.  The
key estimate is independent of any limiting profile and applies to both
orientations of a chain.

For $1\le m\le n$, let
\[
        E_{n,m}=\{n-m+1,\dots,n\}
\]
be the last $m$ labels.  Let $R_{n,m}^+$ and $R_{n,m}^-$ be the lengths of
their longest increasing and decreasing clock subsequences, respectively,
and write $R_{n,m}=R_{n,m}^+$.  For $k\ge0$, define the dyadic terminal block
\begin{equation}\label{eq:dyadic-block}
 D_{n,k}
 =\{i:2^k\le n+1-i<2^{k+1}\}.
\end{equation}

\begin{definition}[Dyadic endpoint comparability]
\label{def:dyadic-comparability}
A triangular array of positive rates has dyadic endpoint constant $K$ if
for every $n$ and every nonempty block $D_{n,k}$,
\begin{equation}\label{eq:dyadic-comparability}
 \frac{\max_{i\in D_{n,k}}\theta_{n,i}}
      {\min_{i\in D_{n,k}}\theta_{n,i}}
 \le K.
\end{equation}
\end{definition}

\begin{lemma}[Dyadic corner-chain estimate]
\label{lem:dyadic-corner}
Assume \eqref{eq:dyadic-comparability}.  There are constants
$a_K,b_K,C_{p,K}>0$ such that, for every $n,m$ and either sign,
\begin{align}
 \Ee e^{b_KR_{n,m}^{\pm}}&\le e^{a_K\sqrt m},
 \label{eq:corner-mgf}\\
 \Pp(R_{n,m}^{\pm}\ge u\sqrt n)
 &\le\exp\{a_K\sqrt m-b_Ku\sqrt n\},
 \qquad u\ge0,
 \label{eq:corner-tail}\\
 \norm{R_{n,m}^{\pm}}_{L^p}&\le C_{p,K}\sqrt m,
 \qquad 1\le p<\infty.
 \label{eq:corner-Lp}
\end{align}
One may take $b_K=1$ and
$a_K=2(2+\sqrt2)\sqrt{eK}$.
\end{lemma}

\begin{proof}
Fix a nonempty subset $B$ of one dyadic block, containing $r$ labels.
Divide its rates by their minimum.  This preserves their relative clock
order, and the normalized rates lie in $[1,K]$.  The increasing
compactification $z=1-e^{-t}$ transforms an exponential clock of rate
$\lambda\in[1,K]$ into a variable with density
\[
        \lambda(1-z)^{\lambda-1}\le K,
        \qquad 0<z<1.
\]
Write $L_B^{\pm}$ for the two chain lengths within $B$.
Lemma~\ref{lem:first-moment-chain} gives
\begin{equation}\label{eq:block-first-moment}
        \Pp(L_B^{\pm}\ge s)\le\frac{(Kr)^s}{(s!)^2},
        \qquad s\ge1.
\end{equation}
For the decreasing orientation the proof is identical: the simplex of
strictly decreasing heights also has volume $1/s!$.
For every $b>0$, the tail-sum formula therefore gives
\begin{align*}
 \Ee e^{bL_B^{\pm}}
 &=1+\sum_{s\ge1}
       (e^{bs}-e^{b(s-1)})\Pp(L_B^{\pm}\ge s)\\
 &\le\sum_{s\ge0}\frac{(Kre^b)^s}{(s!)^2}
 \le e^{2\sqrt{Kre^b}}.
\end{align*}
The last inequality follows from
$\binom{2s}{s}\le4^s$, by comparison with the even terms of the exponential
series.

Put $B_k=D_{n,k}\cap E_{n,m}$ and omit empty intersections.  For either
fixed sign the variables $L_{B_k}^{\pm}$ are independent, and every chain
in $E_{n,m}$ splits into chains in these sets.  Hence
\[
        R_{n,m}^{\pm}\le\sum_k L_{B_k}^{\pm}.
\]
If $J=\lfloor\log_2m\rfloor$, then $|B_k|\le2^k$ and only
$0\le k\le J$ occur.  Consequently
\[
 \sum_k\sqrt{|B_k|}
 \le\sum_{k=0}^J2^{k/2}
 \le(2+\sqrt2)\sqrt m.
\]
Multiplying the exponential-moment bounds with $b=1$ proves
\eqref{eq:corner-mgf} with the stated constants.  Markov's inequality gives
\eqref{eq:corner-tail}; integrating its minimum with $1$ proves
\eqref{eq:corner-Lp}, since $m\ge1$.
\end{proof}

Let $f\in C[0,1]$ satisfy
\begin{equation}\label{eq:endpoint-profile}
        f(x)>0\quad(0\le x<1),
        \qquad f(1)=0.
\end{equation}
Define $A_f,B_f,\tau_f$ and $\rho_f$ by
\eqref{eq:A-B-def}--\eqref{eq:rho-def-intro}.  These formulas remain valid:
$f>0$ almost everywhere implies $A_f(t)\uparrow1$, and $\rho_f$ has uniform
marginals on the open square.  Set
\[
        \rho_f^+(x,y)=\rho_f(x,y),
        \qquad
        \rho_f^-(x,y)=\rho_f(x,1-y).
\]
The reflected density converts decreasing chains into increasing chains.
For these possibly singular densities, endpoint actions are interpreted by
monotone truncation in clock time, as in the following identity.

For $0<\delta<1$, put
\begin{equation}\label{eq:truncated-profile}
        f_\delta(s)=f((1-\delta)s),
        \qquad 0\le s\le1,
\end{equation}
and define
\begin{equation}\label{eq:ell-delta}
        \ell_f^{\pm}(\delta)
        =\sqrt{1-\delta}\,\calL_\pm(f_\delta),
        \qquad \ell_f(\delta)=\ell_f^+(\delta).
\end{equation}
Then
\begin{equation}\label{eq:restriction-identity}
        \ell_f^{\pm}(\delta)
        =\calU_{\rho_f^{\pm}}(1-\delta,1).
\end{equation}
To justify this identity, first restrict the raw clock coordinate to
$0\le t\le T<\infty$.  On this interval the transformation $y=A_f(t)$
and its inverse are continuously differentiable, with strictly positive
derivatives.  The density transformation gives equality of the raw-clock
action and the rank-coordinate action.  Rescaling
$x=(1-\delta)u$ gives the factor $\sqrt{1-\delta}$ on the left.
For the decreasing orientation, reflect the vertical coordinate on each
side of the same change of variables.  The finite-cutoff rank domain is
$y\le A_f(T)$ in the increasing orientation and $y\ge1-A_f(T)$ after
reflection; zero-action axis-parallel segments complete the paths to the
specified endpoints.

Finally let $T\to\infty$.  On $0\le x\le1-\delta$, the profile has a
positive lower bound, so a compactification $z=1-e^{-c_0t}$ with
$0<c_0<\frac12\min_{x\le1-\delta}f(x)$ gives a bounded continuous density.
The omitted tail has vertical width $e^{-c_0T}$, and
Lemma~\ref{lem:U-basic} bounds its possible chain action by
$O(e^{-c_0T/2})$ in either orientation.  Thus the finite-cutoff actions
increase to $\sqrt{1-\delta}\,\calL_\pm(f_\delta)$.
On the rank side, the nonnegative truncated actions increase to the
endpoint actions; taking the supremum over paths commutes with this
increasing limit.  This proves \eqref{eq:restriction-identity} without
requiring a globally Lipschitz inverse at the singular endpoint.

\begin{theorem}[Endpoint-degenerate Luce subsequence laws]
\label{thm:endpoint-luce}
Suppose there are $a_n>0$ and a profile $f$ satisfying
\eqref{eq:endpoint-profile} such that, for every $0<\delta<1$,
\begin{equation}\label{eq:interior-profile-conv}
 \max_{1\le i\le(1-\delta)n}
 \abs{\frac{\theta_{n,i}}{a_n}-f(i/n)}
 \longrightarrow0.
\end{equation}
Assume also dyadic endpoint comparability
\eqref{eq:dyadic-comparability}.  Then the monotone limits
\begin{equation}\label{eq:singular-L-def}
        \calL_\pm^{\mathrm{sing}}(f)
        :=\lim_{\delta\downarrow0}\ell_f^{\pm}(\delta)
        =\lim_{\delta\downarrow0}
          \calU_{\rho_f^{\pm}}(1-\delta,1)
\end{equation}
exist and are finite.  For $0<\delta<1$ they satisfy
\begin{equation}\label{eq:singular-truncation-rate}
 0\le\calL_\pm^{\mathrm{sing}}(f)-\ell_f^{\pm}(\delta)
 \le\frac{a_K}{b_K}\sqrt\delta.
\end{equation}
Moreover,
\begin{align}
        \frac{\LIS(\sigma_n)}{\sqrt n}
        &\longrightarrow\calL_+^{\mathrm{sing}}(f),
        \label{eq:endpoint-LIS-law}\\
        \frac{\LDS(\sigma_n)}{\sqrt n}
        &\longrightarrow\calL_-^{\mathrm{sing}}(f)
        \label{eq:endpoint-LDS-law}
\end{align}
in $L^p$ for every $1\le p<\infty$.  For every $\eps>0$ there are
$c_\eps,C_\eps>0$ such that
\begin{equation}\label{eq:endpoint-concentration}
 \Pp\left(
 \max\left\{
 \abs{\frac{\LIS(\sigma_n)}{\sqrt n}
       -\calL_+^{\mathrm{sing}}(f)},
 \abs{\frac{\LDS(\sigma_n)}{\sqrt n}
       -\calL_-^{\mathrm{sing}}(f)}
 \right\}>\eps
 \right)
 \le C_\eps e^{-c_\eps\sqrt n}
\end{equation}
for all sufficiently large $n$.  Both convergences consequently hold
almost surely under arbitrary couplings.
\end{theorem}

\begin{proof}
Write $L_n^+=\LIS(\sigma_n)$ and $L_n^-=\LDS(\sigma_n)$.
Let $N_{n,\delta}=\lfloor(1-\delta)n\rfloor$, and let
$L_n^{(\delta),\pm}$ be the two chain lengths restricted to labels
$1,\dots,N_{n,\delta}$.  Reindexing these labels by $i/N_{n,\delta}$
changes the profile arguments by $O(n^{-1})$.  Uniform continuity of $f$,
condition \eqref{eq:interior-profile-conv}, and Theorems
\ref{thm:main-luce} and \ref{thm:lds} therefore give
\begin{equation}\label{eq:prefix-limit}
        \frac{L_n^{(\delta),\pm}}{\sqrt n}
        \longrightarrow\ell_f^{\pm}(\delta)
\end{equation}
in $L^p$ for every $1\le p<\infty$, with fixed-deviation probability
$O(e^{-c\sqrt n})$, for each fixed $\delta$.

If $0<\delta'<\delta$, then
$L_n^{(\delta'),\pm}\ge L_n^{(\delta),\pm}$.  Passing to the limits
shows that $\ell_f^{\pm}(\delta)$ is nondecreasing as
$\delta\downarrow0$.  Jensen's inequality and
\eqref{eq:corner-mgf} with $m=n$ give
\[
        \Ee L_n^{\pm}\le\frac{a_K}{b_K}\sqrt n.
\]
Hence both increasing limits in \eqref{eq:singular-L-def} are finite.

The number of labels outside the prefix is
\[
        m_{n,\delta}=n-N_{n,\delta}=\lceil\delta n\rceil.
\]
Deleting these labels from a longest chain gives, for either sign,
\begin{equation}\label{eq:endpoint-truncation-comparison}
 0\le L_n^{\pm}-L_n^{(\delta),\pm}
 \le R_{n,m_{n,\delta}}^{\pm}.
\end{equation}
The same upper bound holds for
$L_n^{(\delta'),\pm}-L_n^{(\delta),\pm}$ when $0<\delta'<\delta$.
Taking expectations, using \eqref{eq:corner-mgf} and Jensen's inequality,
and then letting $n\to\infty$, gives
\[
 0\le\ell_f^{\pm}(\delta')-\ell_f^{\pm}(\delta)
 \le\frac{a_K}{b_K}\sqrt\delta.
\]
Letting $\delta'\downarrow0$ proves
\eqref{eq:singular-truncation-rate}.

For $1\le p<\infty$, \eqref{eq:corner-Lp} and
\eqref{eq:endpoint-truncation-comparison} imply
\[
 \norm{\frac{L_n^{\pm}-L_n^{(\delta),\pm}}{\sqrt n}}_{L^p}
 \le C_{p,K}\sqrt{\delta+n^{-1}}.
\]
The triangle inequality, \eqref{eq:prefix-limit}, and
\eqref{eq:singular-truncation-rate} now prove
\eqref{eq:endpoint-LIS-law}--\eqref{eq:endpoint-LDS-law}, by taking
$n\to\infty$ before $\delta\downarrow0$.

For fixed $\eps>0$, choose $\delta$ so small that
$(a_K/b_K)\sqrt\delta<\eps/3$ and
$a_K\sqrt\delta<b_K\eps/6$.  By \eqref{eq:corner-tail},
\[
 \Pp\left(R_{n,m_{n,\delta}}^{\pm}>
                  \frac{\eps}{3}\sqrt n\right)
 \le
 \exp\left\{a_K\sqrt{\lceil\delta n\rceil}
              -\frac{b_K\eps}{3}\sqrt n\right\}
 \le e^{a_K}e^{-b_K\eps\sqrt n/6}.
\]
Together with the prefix fixed-deviation estimate, the deterministic
comparison, and \eqref{eq:singular-truncation-rate}, this proves
\eqref{eq:endpoint-concentration}, after a union bound over the two signs.
The resulting probabilities are summable in $n$, so Borel--Cantelli gives
almost sure convergence under every coupling.  Finally,
\eqref{eq:restriction-identity} identifies the variational limits in
\eqref{eq:singular-L-def}.
\end{proof}

\subsection{Regularly varying endpoints}

\begin{corollary}[Sampled regularly varying profiles]
\label{cor:regularly-varying-endpoint}
Suppose $f$ satisfies \eqref{eq:endpoint-profile} and
\begin{equation}\label{eq:regular-variation}
        f(1-u)\sim u^\gamma L(1/u),
        \qquad u\downarrow0,
\end{equation}
where $\gamma>0$ and $L$ is slowly varying.  If
\begin{equation}\label{eq:rv-sampling}
 \theta_{n,i}
 =a_n f\left(\frac{i}{n+1}\right)(1+\eps_{n,i}),
 \qquad
 \max_i|\eps_{n,i}|\longrightarrow0,
\end{equation}
then Theorem~\ref{thm:endpoint-luce} applies to both increasing and
decreasing subsequences.
\end{corollary}

\begin{proof}
Interior convergence follows from uniform continuity of $f$.
In a dyadic terminal block, $u=(n+1-i)/(n+1)$ varies by a factor at most
two.  The uniform convergence theorem for regularly varying functions gives
a uniform bound for the ratio of $f(1-u)$ at any two such points once $u$
is small \cite[Theorems~1.2.1 and 1.5.6]{BinghamGoldieTeugels1987}.
A block not contained in this endpoint neighbourhood has all its $u$-values
in a fixed compact subinterval of $(0,1]$, where continuity and positivity
give the same bound.  The relative errors in \eqref{eq:rv-sampling}
change the ratio by a bounded factor for all sufficiently large $n$.
Enlarging the constant to include the finitely many remaining arrays proves
Definition~\ref{def:dyadic-comparability}.
\end{proof}

\begin{corollary}[Both Sukhatme orientations]
\label{cor:sukhatme-orientations}
Let $f_\downarrow(x)=1-x$, and set
\begin{align}
 \calL_{\mathrm{Suk}}^\downarrow
 &=\calL_+^{\mathrm{sing}}(f_\downarrow)
   =\lim_{\delta\downarrow0}
      \calU_{\rho_{\mathrm{Suk}}}(1-\delta,1),
 \label{eq:sukhatme-down-constant}\\
 \calL_{\mathrm{Suk}}^\uparrow
 &=\calL_-^{\mathrm{sing}}(f_\downarrow)
   =\lim_{\delta\downarrow0}
      \calU_{\rho_{\mathrm{Suk}}^-}(1-\delta,1),
 \label{eq:sukhatme-up-constant}
\end{align}
where $\rho_{\mathrm{Suk}}$ is given by
\eqref{eq:suk-rho-intro} and
$\rho_{\mathrm{Suk}}^-(x,y)=\rho_{\mathrm{Suk}}(x,1-y)$.
For Luce permutations $\sigma_n^\downarrow$ with rates $n-i+1$ and
$\sigma_n^\uparrow$ with rates $i$,
\begin{align*}
 \frac1{\sqrt n}
 \bigl(\LIS(\sigma_n^\downarrow),\LDS(\sigma_n^\downarrow)\bigr)
 &\longrightarrow
 \bigl(\calL_{\mathrm{Suk}}^\downarrow,
       \calL_{\mathrm{Suk}}^\uparrow\bigr),\\
 \frac1{\sqrt n}
 \bigl(\LIS(\sigma_n^\uparrow),\LDS(\sigma_n^\uparrow)\bigr)
 &\longrightarrow
 \bigl(\calL_{\mathrm{Suk}}^\uparrow,
       \calL_{\mathrm{Suk}}^\downarrow\bigr).
\end{align*}
These convergences hold in $L^p$ for every $1\le p<\infty$, with the fixed-deviation
bound \eqref{eq:endpoint-concentration}, and almost surely under arbitrary
couplings.  The constant $\calL_{\mathrm{Suk}}$ in
Corollary~\ref{cor:intro-sukhatme} is
$\calL_{\mathrm{Suk}}^\downarrow$.
\end{corollary}

\begin{proof}
Taking $f(x)=1-x$ and $a_n=n+1$ in \eqref{eq:rv-sampling} gives exactly
$\theta_{n,i}=n-i+1$; its dyadic endpoint constant is at most $2$.
Theorem~\ref{thm:endpoint-luce} proves the first convergence and
Corollary~\ref{cor:intro-sukhatme}.

Let $r_n(i)=n+1-i$.  Couple the increasing-rate clocks to the decreasing-rate
clocks by $T_{n,i}^\uparrow=T_{n,r_n(i)}^\downarrow$.
Then $\sigma_n^\uparrow=\sigma_n^\downarrow\circ r_n$.
Reversal of the labels exchanges increasing and decreasing subsequences,
which proves the second convergence with its stated constants and modes.
\end{proof}

\begin{proposition}[Sukhatme truncation bounds]
\label{prop:sukhatme-truncation}
Let $f_\downarrow(x)=1-x$ and
$C_\pm=\calL_\pm^{\mathrm{sing}}(f_\downarrow)$, so that
$C_+=\calL_{\mathrm{Suk}}^\downarrow$ and
$C_-=\calL_{\mathrm{Suk}}^\uparrow$.
For every $0<\delta<1$,
\begin{equation}\label{eq:sukhatme-selfsimilar-bound}
 \ell_{f_\downarrow}^{\pm}(\delta)
 \le C_\pm
 \le\frac{\ell_{f_\downarrow}^{\pm}(\delta)}{1-\sqrt\delta}.
\end{equation}
\end{proposition}

\begin{proof}
Consider the rates $n,n-1,\dots,1$, and put
$N=\lfloor(1-\delta)n\rfloor$ and $m=n-N=\lceil\delta n\rceil$.
The final $m$ clocks have exactly the rates $m,m-1,\dots,1$.
For either orientation, write $L_n^\pm$ for the full chain length,
$P_{n,\delta}^\pm$ for the prefix chain length, and $S_{n,m}^\pm$
for the suffix chain length.  Deterministically,
\[
        P_{n,\delta}^\pm\le L_n^\pm
        \le P_{n,\delta}^\pm+S_{n,m}^\pm,
\]
and $S_{n,m}^\pm$ has the same distribution as $L_m^\pm$.
Take expectations and divide by $\sqrt n$.
The prefix limit \eqref{eq:prefix-limit} and the $L^1$ convergence in
Theorem~\ref{thm:endpoint-luce} give
\[
 \ell_{f_\downarrow}^{\pm}(\delta)
 \le C_\pm
 \le\ell_{f_\downarrow}^{\pm}(\delta)+\sqrt\delta\,C_\pm.
\]
Rearrangement proves \eqref{eq:sukhatme-selfsimilar-bound}.
\end{proof}

\begin{proposition}[An explicit Sukhatme LIS lower bound]
\label{prop:sukhatme-explicit-lower}
The constant for decreasing Sukhatme rates satisfies
\begin{equation}\label{eq:sukhatme-explicit-lower}
        \calL_{\mathrm{Suk}}^\downarrow>3.0602.
\end{equation}
More generally, for every $a>0$,
\begin{equation}\label{eq:sukhatme-trial-action}
 \calL_{\mathrm{Suk}}^\downarrow
 \ge J(a):=4\sqrt a\,e^{-a/2}
                  \int_0^1e^{au^2/2}\,\dd u.
\end{equation}
\end{proposition}

\begin{proof}
For the raw-clock density $(1-x)e^{-(1-x)t}$, consider the increasing curve
\[
        t_a(x)=\frac{ax}{1-x},\qquad 0\le x<1.
\]
Restrict it to $x\le1-\delta$, transform to rank coordinates, and append
a zero-action vertical segment to rank $1$.  This gives an admissible path
for $\ell_{f_\downarrow}^+(\delta)$.  Since
$t_a'(x)=a/(1-x)^2$, its actions increase as $\delta\downarrow0$ to
\[
 2\sqrt a\int_0^1\frac{e^{-ax/2}}{\sqrt{1-x}}\,\dd x
 =4\sqrt a\,e^{-a/2}\int_0^1e^{au^2/2}\,\dd u,
\]
where the equality uses $u=\sqrt{1-x}$.
This proves \eqref{eq:sukhatme-trial-action}.

Take $a=7/4$ and put
\[
 A_{13}=\sum_{j=0}^{13}\frac{(-7/8)^j}{j!},
 \qquad
 B_{12}=\sum_{k=0}^{12}\frac{(7/8)^k}{k!(2k+1)}.
\]
Taylor's formula and positivity of the omitted integral terms give
\[
 0<A_{13}<e^{-7/8},\qquad
 B_{12}<\int_0^1e^{7u^2/8}\,\dd u.
\]
Squaring verifies $\sqrt7>26457513/10^7$, and direct rational arithmetic
gives
\[
 J(7/4)
 >\frac{2\cdot26457513}{10^7}\,A_{13}B_{12}
 >\frac{15301}{5000}=3.0602.
\]
Together with \eqref{eq:sukhatme-trial-action}, this proves
\eqref{eq:sukhatme-explicit-lower}.
\end{proof}

The singularity has a universal local form on comparable endpoint scales.
The next statement is independent of the probabilistic truncation argument.

\begin{proposition}[Regularly varying corner kernel]
\label{prop:corner-kernel}
In addition to \eqref{eq:regular-variation}, suppose
$r(u)=f(1-u)$ is eventually $C^1$, strictly increasing, and
$ur'(u)/r(u)\to\gamma$.  Put
\[
        G_\gamma=\Gamma(1+1/\gamma),
        \qquad
        \psi_\gamma(s)
        =\gamma G_\gamma^\gamma s^\gamma
         e^{-G_\gamma^\gamma s^\gamma}.
\]
Then, locally uniformly for $s\in(0,\infty)$,
\begin{equation}\label{eq:corner-kernel-asymptotic}
        v\,\rho_f(1-sv,1-v)
        \longrightarrow\psi_\gamma(s)
        \qquad(v\downarrow0).
\end{equation}
\end{proposition}

\begin{proof}
Let
$S(t)=1-A_f(t)=\int_0^1e^{-tf(x)}\,\dd x$ and let $h(t)$ be the asymptotic
inverse scale determined by $r(h(t))\sim1/t$.  Karamata's Laplace theorem and
the monotone density theorem give
\[
 S(t)\sim G_\gamma h(t),
 \qquad
 B_f(t)=-S'(t)\sim\frac{S(t)}{\gamma t};
\]
see \cite[Sections~1.5 and 4.12]{BinghamGoldieTeugels1987}.  If
$v=S(t)$ and $u=sv$, regular variation yields
\[
        tr(u)\longrightarrow G_\gamma^\gamma s^\gamma,
\]
locally uniformly for $s\in(0,\infty)$.
Since
$\rho_f(1-u,A_f(t))=r(u)e^{-tr(u)}/B_f(t)$, substitution gives
\eqref{eq:corner-kernel-asymptotic}.
\end{proof}

For the Sukhatme profile $\gamma=1$, this reads
\[
        v\,\rho_{\mathrm{Suk}}(1-sv,1-v)\longrightarrow se^{-s}.
\]

\section{Further questions}
\label{sec:further}

The first-order and localization laws leave several natural problems.

\begin{conjecture}[KPZ fluctuations with exact centering]
\label{conj:TW}
Assume an exactly sampled $C^3$ positive profile and a unique smooth
nondegenerate maximizing path separated from the boundary.  Then there is
$s_f>0$ such that, writing $\chi_2$ for a random variable with the GUE
Tracy--Widom law,
\begin{equation}\label{eq:TW-conjecture}
        \frac{\LIS(\sigma_n)-\Ee\LIS(\sigma_n)}
             {s_fn^{1/6}}
        \Rightarrow \chi_2-\Ee\chi_2.
\end{equation}
Equivalently, one expects
\[
 \frac{\LIS(\sigma_n)-\calL_+(f)\sqrt n}
      {s_fn^{1/6}}
 \Rightarrow\chi_2
\]
whenever the exact finite-$n$ model has no additional deterministic
$n^{1/6}$-order centering term.  In that case
\[
 \Ee\LIS(\sigma_n)
 =\calL_+(f)\sqrt n+s_f(\Ee\chi_2)n^{1/6}+o(n^{1/6}).
\]
\end{conjecture}

The exact-centering qualification is essential.  The first-order hypothesis
$\max_i|\eps_{n,i}|\to0$ allows deterministic profile errors large enough to
shift the LIS by more than $n^{1/6}$.  Centering by the expectation avoids
that ambiguity, but necessarily centers the limiting Tracy--Widom variable as
well.  The uniform case is the Baik--Deift--Johansson theorem
\cite{BaikDeiftJohansson1999}.  Different limits may occur at caustics, for
boundary-touching maximizers, or when the maximizing path is not unique.

\begin{conjecture}[Row-stratified Greene--RSK surface]
\label{conj:RSK-surface}
For a continuous row density $q$, the complete rescaled Young diagram of the
row-stratified word has a deterministic variational limit.  Its Greene
invariants are described by the corresponding optimization over disjoint
families of monotone paths.  Under the rank-coordinate map, this yields a full
RSK limit surface for independent-score and regular Luce rankings.
\end{conjecture}

This would extend the locally uniform permutation theory of Sj\"ostrand
\cite{Sjostrand2023} and the permuton-based work of Dubach
\cite{Dubach2023,Dubach2024} to deterministic one-point-per-fibre arrays.

\begin{conjecture}[Rapidly varying endpoint transition]
\label{conj:rapid-endpoint}
For geometric Sukhatme weights
\[
        \theta_{n,i}=\alpha^{n-i+1},\qquad \alpha>1,
\]
there is a constant $v(\alpha)\in(0,1)$ such that
\[
        \frac{\LIS(\sigma_n)}n\longrightarrow v(\alpha).
\]
More generally, the LIS exponent for endpoint arrays outside the dyadically
comparable class is determined by the number of labels occupying each
multiplicative rate scale.
\end{conjecture}

The geometric model is the simplest regime in which each successive block
contains only $O(1)$ labels per rate decade.  It should therefore lie in a
linear rather than a Hammersley $\sqrt n$ class.  Intermediate arrays, such as
rates varying like $\exp\{\beta_n i/n\}$ with $\beta_n\to\infty$, may exhibit
a crossover governed by the effective number of labels per logarithmic rate
scale.

\begin{conjecture}[Rough row profiles]
\label{conj:rough-q}
Theorem~\ref{thm:stratified-shape} remains valid for bounded measurable row
densities $q$ under a suitable Riemann-approximability hypothesis, with the
upper semicontinuous relaxation of the Hammersley action when $q$ has genuine
discontinuities.
\end{conjecture}

Continuity is used in the proof only to make the within-block oscillation
vanish uniformly and to identify the block functional.  This suggests
extensions to regulated and piecewise continuous profiles and, with a refined
exceptional-block argument, to $L^1$ approximation.

\begin{conjecture}[Moderate deviations and transversal fluctuations]
\label{conj:moderate}
In the unique interior-geodesic regime, the fixed-deviation estimate of this
paper admits a uniform moderate-deviation interpolation down to the
$n^{1/6}$ longitudinal and $n^{-1/6}$ transversal scales, with the usual KPZ
exponents and profile-dependent curvature constants.
\end{conjecture}

\begin{remark}[Scope of the proved concentration]
The bounds proved here concern fixed deviations on the $\sqrt n$ first-order
scale.  They do not by themselves imply Tracy--Widom fluctuations, moderate
deviations, or KPZ transversal exponents.
\end{remark}


\begin{thebibliography}{99}

\bibitem{AldousDiaconis1995}
D. Aldous and P. Diaconis.
\newblock Hammersley's interacting particle process and longest increasing
subsequences.
\newblock \emph{Probab. Theory Related Fields} 103 (1995), no.~2, 199--213.
\newblock \href{https://doi.org/10.1007/BF01204214}{doi:10.1007/BF01204214}.

\bibitem{BaikDeiftJohansson1999}
J. Baik, P. Deift, and K. Johansson.
\newblock On the distribution of the length of the longest increasing
subsequence of random permutations.
\newblock \emph{J. Amer. Math. Soc.} 12 (1999), no.~4, 1119--1178.
\newblock \href{https://doi.org/10.1090/S0894-0347-99-00307-0}
{doi:10.1090/S0894-0347-99-00307-0}.

\bibitem{BhatnagarPeled2015}
N. Bhatnagar and R. Peled.
\newblock Lengths of monotone subsequences in a Mallows permutation.
\newblock \emph{Probab. Theory Related Fields} 161 (2015), nos.~3--4,
719--780.
\newblock \href{https://doi.org/10.1007/s00440-014-0559-7}
{doi:10.1007/s00440-014-0559-7}.

\bibitem{BinghamGoldieTeugels1987}
N. H. Bingham, C. M. Goldie, and J. L. Teugels.
\newblock \emph{Regular Variation}.
\newblock Encyclopedia of Mathematics and its Applications, vol.~27,
Cambridge University Press, Cambridge, 1987.
\newblock \href{https://doi.org/10.1017/CBO9780511721434}
{doi:10.1017/CBO9780511721434}.

\bibitem{BorgaChatterjeeDiaconis2026}
J. Borga, S. Chatterjee, and P. Diaconis.
\newblock Permuton and local limits for the Luce model.
\newblock arXiv:2509.07729v3 (2026), accepted for publication in
\emph{Ann. Appl. Probab.}
\newblock \href{https://doi.org/10.48550/arXiv.2509.07729}
{doi:10.48550/arXiv.2509.07729}.

\bibitem{CalderEsedogluHero2014}
J. Calder, S. Esedo\u{g}lu, and A. O. Hero.
\newblock A Hamilton--Jacobi equation for the continuum limit of
nondominated sorting.
\newblock \emph{SIAM J. Math. Anal.} 46 (2014), no.~1, 603--638.
\newblock \href{https://doi.org/10.1137/13092842X}
{doi:10.1137/13092842X}.

\bibitem{ChatterjeeDiaconisKim2024}
S. Chatterjee, P. Diaconis, and G. B. Kim.
\newblock Enumerative theory for the Tsetlin library.
\newblock \emph{J. Algebra} 655 (2024), 139--162.
\newblock \href{https://doi.org/10.1016/j.jalgebra.2023.08.009}
{doi:10.1016/j.jalgebra.2023.08.009}.

\bibitem{DeuschelZeitouni1995}
J.-D. Deuschel and O. Zeitouni.
\newblock Limiting curves for i.i.d. records.
\newblock \emph{Ann. Probab.} 23 (1995), no.~2, 852--878.
\newblock \href{https://doi.org/10.1214/aop/1176988293}
{doi:10.1214/aop/1176988293}.

\bibitem{DeuschelZeitouni1999}
J.-D. Deuschel and O. Zeitouni.
\newblock On increasing subsequences of i.i.d. samples.
\newblock \emph{Combin. Probab. Comput.} 8 (1999), no.~3, 247--263.
\newblock \href{https://doi.org/10.1017/S0963548399003776}
{doi:10.1017/S0963548399003776}.

\bibitem{Dubach2023}
V. Dubach.
\newblock Locally uniform random permutations with large increasing
subsequences.
\newblock \emph{Combinatorial Theory} 3 (2023), no.~3, Paper No.~8.
\newblock \href{https://doi.org/10.5070/C63362784}
{doi:10.5070/C63362784}.

\bibitem{Dubach2024}
V. Dubach.
\newblock Increasing subsequences of linear size in random permutations and
the Robinson--Schensted tableaux of permutons.
\newblock \emph{Random Structures Algorithms} 65 (2024), no.~3, 488--534.
\newblock \href{https://doi.org/10.1002/rsa.21223}
{doi:10.1002/rsa.21223}.

\bibitem{ErdosSzekeres1935}
P. Erd\H{o}s and G. Szekeres.
\newblock A combinatorial problem in geometry.
\newblock \emph{Compositio Math.} 2 (1935), 463--470.

\bibitem{LoganShepp1977}
B. F. Logan and L. A. Shepp.
\newblock A variational problem for random Young tableaux.
\newblock \emph{Adv. Math.} 26 (1977), no.~2, 206--222.
\newblock \href{https://doi.org/10.1016/0001-8708(77)90030-5}
{doi:10.1016/0001-8708(77)90030-5}.

\bibitem{Luce1959}
R. D. Luce.
\newblock \emph{Individual Choice Behavior: A Theoretical Analysis}.
\newblock Wiley, New York, 1959.

\bibitem{Mosteller1951}
F. Mosteller.
\newblock Remarks on the method of paired comparisons: I. The least squares
solution assuming equal standard deviations and equal correlations.
\newblock \emph{Psychometrika} 16 (1951), no.~1, 3--9.
\newblock \href{https://doi.org/10.1007/BF02313422}
{doi:10.1007/BF02313422}.

\bibitem{MuellerStarr2013}
C. Mueller and S. Starr.
\newblock The length of the longest increasing subsequence of a random
Mallows permutation.
\newblock \emph{J. Theoret. Probab.} 26 (2013), no.~2, 514--540.
\newblock \href{https://doi.org/10.1007/s10959-011-0364-5}
{doi:10.1007/s10959-011-0364-5}.

\bibitem{Sjostrand2023}
J. Sj\"ostrand.
\newblock Monotone subsequences in locally uniform random permutations.
\newblock \emph{Ann. Probab.} 51 (2023), no.~4, 1502--1547.
\newblock \href{https://doi.org/10.1214/23-AOP1624}
{doi:10.1214/23-AOP1624}.

\bibitem{Talagrand1995}
M. Talagrand.
\newblock Concentration of measure and isoperimetric inequalities in product
spaces.
\newblock \emph{Publ. Math. Inst. Hautes \'{E}tudes Sci.} 81 (1995), 73--205.
\newblock \href{https://doi.org/10.1007/BF02699376}
{doi:10.1007/BF02699376}.

\bibitem{Thurstone1927}
L. L. Thurstone.
\newblock A law of comparative judgment.
\newblock \emph{Psychological Review} 34 (1927), no.~4, 273--286.

\bibitem{VershikKerov1977}
A. M. Vershik and S. V. Kerov.
\newblock Asymptotics of the Plancherel measure of the symmetric group and the
limiting form of Young tableaux.
\newblock \emph{Dokl. Akad. Nauk SSSR} 233 (1977), no.~6, 1024--1027;
English translation in \emph{Soviet Math. Dokl.} 18 (1977), 527--531.

\end{thebibliography}

\end{document}
